% Options for packages loaded elsewhere
\PassOptionsToPackage{unicode}{hyperref}
\PassOptionsToPackage{hyphens}{url}
%
\documentclass[
]{article}
\usepackage{amsmath,amssymb}
\usepackage{lmodern}
\usepackage{iftex}
\ifPDFTeX
  \usepackage[T1]{fontenc}
  \usepackage[utf8]{inputenc}
  \usepackage{textcomp} % provide euro and other symbols
\else % if luatex or xetex
  \usepackage{unicode-math}
  \defaultfontfeatures{Scale=MatchLowercase}
  \defaultfontfeatures[\rmfamily]{Ligatures=TeX,Scale=1}
\fi
% Use upquote if available, for straight quotes in verbatim environments
\IfFileExists{upquote.sty}{\usepackage{upquote}}{}
\IfFileExists{microtype.sty}{% use microtype if available
  \usepackage[]{microtype}
  \UseMicrotypeSet[protrusion]{basicmath} % disable protrusion for tt fonts
}{}
\makeatletter
\@ifundefined{KOMAClassName}{% if non-KOMA class
  \IfFileExists{parskip.sty}{%
    \usepackage{parskip}
  }{% else
    \setlength{\parindent}{0pt}
    \setlength{\parskip}{6pt plus 2pt minus 1pt}}
}{% if KOMA class
  \KOMAoptions{parskip=half}}
\makeatother
\usepackage{xcolor}
\usepackage[margin=1in]{geometry}
\usepackage{color}
\usepackage{fancyvrb}
\newcommand{\VerbBar}{|}
\newcommand{\VERB}{\Verb[commandchars=\\\{\}]}
\DefineVerbatimEnvironment{Highlighting}{Verbatim}{commandchars=\\\{\}}
% Add ',fontsize=\small' for more characters per line
\usepackage{framed}
\definecolor{shadecolor}{RGB}{248,248,248}
\newenvironment{Shaded}{\begin{snugshade}}{\end{snugshade}}
\newcommand{\AlertTok}[1]{\textcolor[rgb]{0.94,0.16,0.16}{#1}}
\newcommand{\AnnotationTok}[1]{\textcolor[rgb]{0.56,0.35,0.01}{\textbf{\textit{#1}}}}
\newcommand{\AttributeTok}[1]{\textcolor[rgb]{0.77,0.63,0.00}{#1}}
\newcommand{\BaseNTok}[1]{\textcolor[rgb]{0.00,0.00,0.81}{#1}}
\newcommand{\BuiltInTok}[1]{#1}
\newcommand{\CharTok}[1]{\textcolor[rgb]{0.31,0.60,0.02}{#1}}
\newcommand{\CommentTok}[1]{\textcolor[rgb]{0.56,0.35,0.01}{\textit{#1}}}
\newcommand{\CommentVarTok}[1]{\textcolor[rgb]{0.56,0.35,0.01}{\textbf{\textit{#1}}}}
\newcommand{\ConstantTok}[1]{\textcolor[rgb]{0.00,0.00,0.00}{#1}}
\newcommand{\ControlFlowTok}[1]{\textcolor[rgb]{0.13,0.29,0.53}{\textbf{#1}}}
\newcommand{\DataTypeTok}[1]{\textcolor[rgb]{0.13,0.29,0.53}{#1}}
\newcommand{\DecValTok}[1]{\textcolor[rgb]{0.00,0.00,0.81}{#1}}
\newcommand{\DocumentationTok}[1]{\textcolor[rgb]{0.56,0.35,0.01}{\textbf{\textit{#1}}}}
\newcommand{\ErrorTok}[1]{\textcolor[rgb]{0.64,0.00,0.00}{\textbf{#1}}}
\newcommand{\ExtensionTok}[1]{#1}
\newcommand{\FloatTok}[1]{\textcolor[rgb]{0.00,0.00,0.81}{#1}}
\newcommand{\FunctionTok}[1]{\textcolor[rgb]{0.00,0.00,0.00}{#1}}
\newcommand{\ImportTok}[1]{#1}
\newcommand{\InformationTok}[1]{\textcolor[rgb]{0.56,0.35,0.01}{\textbf{\textit{#1}}}}
\newcommand{\KeywordTok}[1]{\textcolor[rgb]{0.13,0.29,0.53}{\textbf{#1}}}
\newcommand{\NormalTok}[1]{#1}
\newcommand{\OperatorTok}[1]{\textcolor[rgb]{0.81,0.36,0.00}{\textbf{#1}}}
\newcommand{\OtherTok}[1]{\textcolor[rgb]{0.56,0.35,0.01}{#1}}
\newcommand{\PreprocessorTok}[1]{\textcolor[rgb]{0.56,0.35,0.01}{\textit{#1}}}
\newcommand{\RegionMarkerTok}[1]{#1}
\newcommand{\SpecialCharTok}[1]{\textcolor[rgb]{0.00,0.00,0.00}{#1}}
\newcommand{\SpecialStringTok}[1]{\textcolor[rgb]{0.31,0.60,0.02}{#1}}
\newcommand{\StringTok}[1]{\textcolor[rgb]{0.31,0.60,0.02}{#1}}
\newcommand{\VariableTok}[1]{\textcolor[rgb]{0.00,0.00,0.00}{#1}}
\newcommand{\VerbatimStringTok}[1]{\textcolor[rgb]{0.31,0.60,0.02}{#1}}
\newcommand{\WarningTok}[1]{\textcolor[rgb]{0.56,0.35,0.01}{\textbf{\textit{#1}}}}
\usepackage{graphicx}
\makeatletter
\def\maxwidth{\ifdim\Gin@nat@width>\linewidth\linewidth\else\Gin@nat@width\fi}
\def\maxheight{\ifdim\Gin@nat@height>\textheight\textheight\else\Gin@nat@height\fi}
\makeatother
% Scale images if necessary, so that they will not overflow the page
% margins by default, and it is still possible to overwrite the defaults
% using explicit options in \includegraphics[width, height, ...]{}
\setkeys{Gin}{width=\maxwidth,height=\maxheight,keepaspectratio}
% Set default figure placement to htbp
\makeatletter
\def\fps@figure{htbp}
\makeatother
\setlength{\emergencystretch}{3em} % prevent overfull lines
\providecommand{\tightlist}{%
  \setlength{\itemsep}{0pt}\setlength{\parskip}{0pt}}
\setcounter{secnumdepth}{-\maxdimen} % remove section numbering
\ifLuaTeX
  \usepackage{selnolig}  % disable illegal ligatures
\fi
\IfFileExists{bookmark.sty}{\usepackage{bookmark}}{\usepackage{hyperref}}
\IfFileExists{xurl.sty}{\usepackage{xurl}}{} % add URL line breaks if available
\urlstyle{same} % disable monospaced font for URLs
\hypersetup{
  hidelinks,
  pdfcreator={LaTeX via pandoc}}

\author{}
\date{\vspace{-2.5em}}

\begin{document}

\hypertarget{description}{%
\section{DESCRIPTION}\label{description}}

\begin{verbatim}
Package: ShedsHT
Title: The SHEDS-HT model for estimating human exposure to chemicals.
Version: 0.1.10
Author: Kristin Isaacs [aut, cre]
Maintainer: Kristin Isaacs <isaacs.kristin@epa.gov>
Description: The ShedsHT R package runs the Stochastic Human Exposure and Dose
    Simulation-High Throughput screening model which estmates human exposure to a
    wide range of chemicals.The people in SHEDS-HT are simulated individuals who
    collectively form a representative sample of the target population, as chosen by
    the user. The model is cross-sectional, with just one simulated day (24 hours)
    for each simulated person, although the selected day is not necessarily the same
    from one person to another. SHEDS-HT is stochastic, which means that many inputs
    are sampled randomly from user-specified distributions that are intended to
    capture variability. In the SHEDS series of models, variability and uncertainty
    are typically handled by a two-stage Monte Carlo process, but SHEDS-HT currently
    has a single stage and does not directly estimate uncertainty.
License: MIT
Encoding: UTF-8
LazyData: true
RoxygenNote: 7.2.1
Imports: data.table, ggplot2, stringr, plyr
Suggests: knitr, rmarkdown
VignetteBuilder: knitr
NeedsCompilation: no
Packaged: 2019-08-27 14:40:53 UTC; kisaacs1
\end{verbatim}

\hypertarget{act.diary.pools-act.diary.pools}{%
\section{\texorpdfstring{\texttt{act.diary.pools}:
act.diary.pools}{act.diary.pools: act.diary.pools}}\label{act.diary.pools-act.diary.pools}}

\hypertarget{description-1}{%
\subsection{Description}\label{description-1}}

Assigns activity diaries from the
\texttt{{[}Activity\_diaries{]}(Activity\_diaries)} input on the
\texttt{{[}Run{]}(Run)} file (read through the
\texttt{{[}read.act.diaries{]}(read.act.diaries)} function) to pools
based on age, gender, and season

\hypertarget{usage}{%
\subsection{Usage}\label{usage}}

\begin{Shaded}
\begin{Highlighting}[]
\FunctionTok{act.diary.pools}\NormalTok{(diaries, specs)}
\end{Highlighting}
\end{Shaded}

\hypertarget{arguments}{%
\subsection{Arguments}\label{arguments}}

\begin{itemize}
\tightlist
\item
  \texttt{diaries}: A data set created internally in SHEDS.HT through
  the \texttt{{[}read.act.diaries{]}(read.act.diaries)} function. The
  data are activity diaries, which indicate the amount of time and level
  of metabolic activity in various micro environments. Each line of data
  represents one person-day (24 hours).
\item
  \texttt{specs}: Output of the
  \texttt{{[}read.run.file{]}(read.run.file)} function, which can be
  modified by the \texttt{{[}update.specs{]}(update.specs)} function
  before input into \texttt{act.diary.pools}.
\end{itemize}

\hypertarget{details}{%
\subsection{Details}\label{details}}

The \texttt{act.diary.pools} function assigns the activity diaries to
pools. A given diary may belong to many pools, since every year of age
has its own pool. Large pools may contain a list of several hundred
diary numbers. The code contains four loops and on each step performs a
sub-setting of the list of diary numbers.

\hypertarget{note}{%
\subsection{Note}\label{note}}

The input to \texttt{act.diary.pools}, diaries, is created by reading in
the Activity\_diaries input file (specified on the
\texttt{{[}Run{]}(Run)} file) with the
\texttt{{[}read.act.diaries{]}(read.act.diaries)} function within the
\texttt{{[}run{]}(run)} function.

\hypertarget{author}{%
\subsection{Author}\label{author}}

Kristin Isaacs, Graham Glen

\hypertarget{seealso}{%
\subsection{Seealso}\label{seealso}}

\texttt{{[}run{]}(run)}

\hypertarget{value}{%
\subsection{Value}\label{value}}

pool A vector of lists. The length of \texttt{pool} is the product of
the genders, seasons, and ages inputs specified in the
\texttt{{[}Run{]}(Run)} file. Each element is a list of acceptable
activity diary numbers for each year of age, gender, weekend, and season
combination. For example, \texttt{pool}{[}100{]} may have the name
M0P99, which indicates that it is for males, on weekdays, in spring,for
age=99. In addition, If the function runs successfully, the following
message will be printed: ``Activity Diary Pooling completed''

\hypertarget{add.factors-add.factors}{%
\section{\texorpdfstring{\texttt{add.factors}:
add.factors}{add.factors: add.factors}}\label{add.factors-add.factors}}

Adds specific exposure factors to the \texttt{pdm} data table, which is
output from the \texttt{{[}add.media{]}(add.media)} function. Here,
``specific'' means taking into account the age, gender, season, and
exposure media for each person.

\hypertarget{description-2}{%
\subsection{Description}\label{description-2}}

add.factors Adds specific exposure factors to the \texttt{pdm} data
table, which is output from the \texttt{{[}add.media{]}(add.media)}
function. Here, ``specific'' means taking into account the age, gender,
season, and exposure media for each person.

\hypertarget{usage-1}{%
\subsection{Usage}\label{usage-1}}

\begin{Shaded}
\begin{Highlighting}[]
\FunctionTok{add.factors}\NormalTok{(n, gen.f, med.f, exp.f, surf, pdm)}
\end{Highlighting}
\end{Shaded}

\hypertarget{arguments-1}{%
\subsection{Arguments}\label{arguments-1}}

\begin{itemize}
\tightlist
\item
  \texttt{n}: Number of persons
\item
  \texttt{gen.f}: Non-media specific exposure factors as a data table.
  Output from the \texttt{{[}gen.factor.tables{]}(gen.factor.tables)}
  function.
\item
  \texttt{med.f}: Media specific exposure factors presented as a data
  table. Output from the
  \texttt{{[}med.factor.tables{]}(med.factor.tables)} function.
\item
  \texttt{exp.f}: Distributional parameters for the exposure factors.
  Output of the \texttt{{[}exp.factors{]}(exp.factors)} function.
\item
  \texttt{surf}: A list of surface media. Modified output of the
  \texttt{{[}read.media.file{]}(read.media.file)} function.
\item
  \texttt{pdm}: A data table containing the \texttt{pd} data frame of
  physiological and demographic parameters for each theoretical person,
  and the \texttt{dur} array, which specifies the duration of exposure
  to each potential exposure medium for each person in \texttt{pd}.
  Output of the \texttt{{[}add.media{]}(add.media)} function.
\end{itemize}

\hypertarget{details-1}{%
\subsection{Details}\label{details-1}}

The process of adding specific exposure factors to \texttt{pdm} involves
multiple steps. First, \texttt{w} is determined, which is the number of
general factors plus the product of the number of media-specific factors
and the number of surface media. Air media do not have media specific
factors in this version of SHEDS. An array, \texttt{q}, of uniform
random samples is generated, with one row per person and \texttt{w}
columns. A zero matrix, \texttt{r}, of the same size is defined. Once
these matrices are defined, the media-specific factors are determined.
Two nested loops over variable and surface type generate the values,
which are stored in \texttt{r}. Next, another FOR loop determines the
general factors. The \texttt{p} data set contains the age, gender, and
season for each person. These two data sets are then merged. The
evaluation of these factors is handled by the
\texttt{{[}eval.factors{]}(eval.factors)} function. One of the exposure
factors is \texttt{handwash.freq}. This was also part of
SHEDS-Multimedia, where it represented the mean number of hours in the
day with hand washing events. An important aspect of that model was that
because each person was followed longitudinally, the actual number of
hand washes on each day varied from one day to the next. Because of
this, the distribution for \texttt{handwash.freq} did not need to be
restricted to integer values, as (for example) a mean of 4.5 per day is
acceptable and achievable, while choosing integer numbers of hand washes
each day. One of the early goals with SHEDS.HT was to attempt to
reproduce selected results from SHEDS-Multimedia. Therefore, similar
logic was built into the current model. The \texttt{hand.washes}
variable is sampled from a distribution centered on
\texttt{handwash.freq}, and then rounded to the nearest integer. The
\texttt{bath} variable is another difficult concept. In theory, baths
and showers are recorded on the activity diaries. In practice, the
activity diaries were constructed from approximately 20 separate
studies, some of which did not contain enough detail to identify
separate bath or shower events. The result is that about half of all
diaries record such events, but the true rate in the population is
higher. The \texttt{bath.p} variable was created to address this. It
represents the probability that a non bath/shower activity diary should
actually have one. Therefore, if the diary has one, then SHEDS
automatically has one. Otherwise, a binomial sample using
\texttt{bath.p} as the probability is drawn. A bath/shower occurs unless
both of these are zero. The effectiveness of hand washes or bath/shower
at removing chemical from the skin is determined in the
\texttt{{[}post.exposure{]}(post.exposure)} function.

\hypertarget{author-1}{%
\subsection{Author}\label{author-1}}

Kristin Isaacs, Graham Glen

\hypertarget{seealso-1}{%
\subsection{Seealso}\label{seealso-1}}

\texttt{{[}eval.factors{]}(eval.factors)},
\texttt{{[}post.exposure{]}(post.exposure)}

\hypertarget{value-1}{%
\subsection{Value}\label{value-1}}

pdmf A data set containing the \texttt{pdm} data table as well as media
specific exposure factors, the number of baths taken, and the number of
hand wash events occurring per day per person contained in \texttt{pdm}.

\hypertarget{add.fugs-add.fugs}{%
\section{\texorpdfstring{\texttt{add.fugs}:
add.fugs}{add.fugs: add.fugs}}\label{add.fugs-add.fugs}}

\hypertarget{description-3}{%
\subsection{Description}\label{description-3}}

Evaluates the variables in the fugacity input file and creates values
for those variables corresponding to each simulated person.

\hypertarget{usage-2}{%
\subsection{Usage}\label{usage-2}}

\begin{Shaded}
\begin{Highlighting}[]
\FunctionTok{add.fugs}\NormalTok{(n.per, x, pdmf)}
\end{Highlighting}
\end{Shaded}

\hypertarget{arguments-2}{%
\subsection{Arguments}\label{arguments-2}}

\begin{itemize}
\tightlist
\item
  \texttt{n.per}: The total number of simulated persons in this model
  run specified in the \texttt{{[}Run{]}(Run)} file
\item
  \texttt{x}: the output of the
  \texttt{{[}read.fug.inputs{]}(read.fug.inputs)} function.
\item
  \texttt{pdmf}: he output of the
  \texttt{{[}add.factors{]}(add.factors)} function. A data set
  containing physiological and demographic parameters for each
  theoretical person, the duration of exposure to each potential
  exposure medium for each person, the media specific exposure factors,
  and the number of baths taken and hand wash events occurring per day
  for each person.
\end{itemize}

\hypertarget{keyword}{%
\subsection{Keyword}\label{keyword}}

SHEDS

\hypertarget{details-2}{%
\subsection{Details}\label{details-2}}

This function evaluates the variables in the fugacity input file,
creating one value for each simulated person and adding each variable as
a new column in the pdmf data set (renamed as pdmff). All variables are
left i n their original units except for those with units ug/cm2, which
are converted to ug/m2. In the fugacity calculations, all masses are in
ug and all lengths are in m.

\hypertarget{author-2}{%
\subsection{Author}\label{author-2}}

Kristin Isaacs, Graham Glen

\hypertarget{seealso-2}{%
\subsection{Seealso}\label{seealso-2}}

\texttt{{[}add.factors{]}(add.factors)},
\texttt{{[}add.media{]}(add.media)},
\texttt{{[}select\_people{]}(select\_people)},
\texttt{{[}Fugacity{]}(Fugacity)}, \texttt{{[}run{]}(run)}

\hypertarget{value-2}{%
\subsection{Value}\label{value-2}}

pdmff Output contains values sampled from the distributions of each
relevant variable in the \texttt{{[}Fugacity{]}(Fugacity)} input file
for each theoretical person.

\hypertarget{add.media-add.media}{%
\section{\texorpdfstring{\texttt{add.media}:
add.media}{add.media: add.media}}\label{add.media-add.media}}

\hypertarget{description-4}{%
\subsection{Description}\label{description-4}}

For each theoretical person parameterized in the \texttt{pd} data frame,
which is output from the
\texttt{{[}select.people{]}(select.people)}function, this function
generates the exposure duration for each potential exposure medium.

\hypertarget{usage-3}{%
\subsection{Usage}\label{usage-3}}

\begin{Shaded}
\begin{Highlighting}[]
\FunctionTok{add.media}\NormalTok{(n, media, pd)}
\end{Highlighting}
\end{Shaded}

\hypertarget{arguments-3}{%
\subsection{Arguments}\label{arguments-3}}

\begin{itemize}
\tightlist
\item
  \texttt{n}: Number of persons.
\item
  \texttt{media}: List of the potential contact media for the model;
  output of the \texttt{{[}read.media.file{]}(read.media.file)}
  function. Each is found in a specific microenvironment (micro).
\item
  \texttt{pd}: Data frame containing internally assigned demographic and
  physiological parameters for each theoretical person modeled in
  SHEDS.HT. Output of the \texttt{{[}select.people{]}(select.people)}
  function
\end{itemize}

\hypertarget{details-3}{%
\subsection{Details}\label{details-3}}

To generate the exposure duration for each theoretical person, an array
q of uniform random samples is generated with rows = \texttt{n} (per
set, where set size is specified by the user in the
\texttt{{[}Run{]}(Run)} input file) and columns = \texttt{nrow(media)}
(the number of potential exposure media). An array called \texttt{dur}
is created for the number of minutes on the activity diary in the
relevant micro multiplied by the relevant probability of contact (as
specified in the \texttt{media} input). The array consists of a row for
each person and a column for each of the exposure media. The output is a
data table containing the \texttt{pd} data frame and the \texttt{dur}
array.

\hypertarget{author-3}{%
\subsection{Author}\label{author-3}}

Kristin Isaacs, Graham Glen

\hypertarget{seealso-3}{%
\subsection{Seealso}\label{seealso-3}}

\texttt{{[}select.people{]}(select.people)},
\texttt{{[}read.media.file{]}(read.media.file)}

\hypertarget{value-3}{%
\subsection{Value}\label{value-3}}

pdm A data table containing the \texttt{pd} data frame of physiological
and demographic parameters for each theoretical person, and the
\texttt{dur} array, which specifies the duration of exposure to each
potential exposure medium for each person in \texttt{pd}.

\hypertarget{allstats.variable.rank.plot-allstats.variable.rank.plot}{%
\section{\texorpdfstring{\texttt{allstats.variable.rank.plot}:
allstats.variable.rank.plot}{allstats.variable.rank.plot: allstats.variable.rank.plot}}\label{allstats.variable.rank.plot-allstats.variable.rank.plot}}

\hypertarget{usage-4}{%
\subsection{Usage}\label{usage-4}}

\begin{Shaded}
\begin{Highlighting}[]
\FunctionTok{allstats.variable.rank.plot}\NormalTok{(}
  \AttributeTok{run.name =}\NormalTok{ specs}\SpecialCharTok{$}\NormalTok{run.name,}
  \AttributeTok{output.variable =} \StringTok{"abs.tot.mgkg"}\NormalTok{,}
  \AttributeTok{metrics =} \FunctionTok{c}\NormalTok{(}\StringTok{"50\%"}\NormalTok{, }\StringTok{"95\%"}\NormalTok{, }\StringTok{"mean"}\NormalTok{),}
  \AttributeTok{cohort.col =} \StringTok{"Total"}
\NormalTok{)}
\end{Highlighting}
\end{Shaded}

\hypertarget{arguments-4}{%
\subsection{Arguments}\label{arguments-4}}

\begin{itemize}
\tightlist
\item
  \texttt{run.name}: default = name of last run (in current R session)
\item
  \texttt{output.variable}: default = abs.tot.mgkg (mg/kg/day).
\item
  \texttt{metrics}: argument is a vector of character strings with
  desired percentiles or metrics. default =
  c(``50\itemcohort.colagurment vector of different cohorts. default
  =''Total''. To add cohorts use a vector of concatenated character
  strings, such as cohort.col = c(``Total'', ``Male'', ``Female'')a plot
  of output.variable vs chemical rankallstats.variable.rank.plot is a
  ggplot scatter plot with output.variable on the y axis and chemical
  rank on the x axis. User can color points by cohort (see
  ``cohort.col'' argument) and assign percentiles to be points (see
  ``metrics'' argument)A SHEDS run creates an output file for each
  chemical. This plot pulls from the allstats.csv for each chemical in
  the output file. Chemicals are ranked in ascending order of
  output.variable on the x axis. The y axis is log transformed and a
  constant of 1e-10 added to y variable for log transformation of 0's.
  The default color is black for the ``Total'' cohort, but user can
  include as many cohorts as they wish as long as they are calculated in
  the SHEDS run. Percentiles must be calculated in SHEDS run.the
  selected cohorts.
\end{itemize}

\hypertarget{check.foods-check.foods}{%
\section{\texorpdfstring{\texttt{check.foods}:
check.foods}{check.foods: check.foods}}\label{check.foods-check.foods}}

This function creates a list of unique food types.

\hypertarget{description-5}{%
\subsection{Description}\label{description-5}}

check.foods This function creates a list of unique food types.

\hypertarget{usage-5}{%
\subsection{Usage}\label{usage-5}}

\begin{Shaded}
\begin{Highlighting}[]
\FunctionTok{check.foods}\NormalTok{(s)}
\end{Highlighting}
\end{Shaded}

\hypertarget{check.src.scen.flags-check.src.scen.flags}{%
\section{\texorpdfstring{\texttt{check.src.scen.flags}:
check.src.scen.flags}{check.src.scen.flags: check.src.scen.flags}}\label{check.src.scen.flags-check.src.scen.flags}}

\hypertarget{description-6}{%
\subsection{Description}\label{description-6}}

This function checks whether the settings on the source.scen object
(input as ``df'') are set to numeric values of 1 or 0. If the column for
a particular exposure scenario is missing from ``df'', it is created and
assigned ``0'' for all sources.

\hypertarget{usage-6}{%
\subsection{Usage}\label{usage-6}}

\begin{Shaded}
\begin{Highlighting}[]
\FunctionTok{check.src.scen.flags}\NormalTok{(df)}
\end{Highlighting}
\end{Shaded}

\hypertarget{check.src.scen.types-check.src.scen.types}{%
\section{\texorpdfstring{\texttt{check.src.scen.types}:
Check.src.scen.types}{check.src.scen.types: Check.src.scen.types}}\label{check.src.scen.types-check.src.scen.types}}

\hypertarget{description-7}{%
\subsection{Description}\label{description-7}}

This function checks whether the settings on the source.scen object
(input as ``df'') are consistent with the source.type for each source.

\hypertarget{usage-7}{%
\subsection{Usage}\label{usage-7}}

\begin{Shaded}
\begin{Highlighting}[]
\FunctionTok{check.src.scen.types}\NormalTok{(df)}
\end{Highlighting}
\end{Shaded}

\hypertarget{chem.fug-chem.fug}{%
\section{\texorpdfstring{\texttt{chem.fug}:
chem.fug}{chem.fug: chem.fug}}\label{chem.fug-chem.fug}}

\hypertarget{description-8}{%
\subsection{Description}\label{description-8}}

Creates distributions of chemical specific parameters for each chemical
of interest, in order to reflect real-world variability and uncertainty.
These distributions are then sampled to create chemical specific
parameters associated with the exposure of each simulated person.

\hypertarget{usage-8}{%
\subsection{Usage}\label{usage-8}}

\begin{Shaded}
\begin{Highlighting}[]
\FunctionTok{chem.fug}\NormalTok{(n.per, cprops, x)}
\end{Highlighting}
\end{Shaded}

\hypertarget{arguments-5}{%
\subsection{Arguments}\label{arguments-5}}

\begin{itemize}
\tightlist
\item
  \texttt{n.per}: The total number of simulated persons in this model
  run specified in the \texttt{{[}Run{]}(Run)} file.
\item
  \texttt{cprops}: The chemical properties required for SHEDS-HT (output
  of \texttt{chem.props} function). The default data was prepared from
  publicly available databases using a custom program (not part of
  SHEDS-HT). The default file contains about 17 numerical inputs per
  chemical, but most are not used. The required properties are molecular
  weight (\texttt{MW}), vapor pressure (\texttt{VP.Pa}), solubility
  (\texttt{water.sol.mg.l}), octanol-water partition coefficient
  (\texttt{log.Kow}), air decay rate (\texttt{half.air.hr}), decay rate
  on surfaces (\texttt{half.sediment.hr}), and permeability coefficient
  (\texttt{Kp}).
\item
  \texttt{x}: From the output of the
  \texttt{{[}read.fug.inputs{]}(read.fug.inputs)} function.
\end{itemize}

\hypertarget{keyword-1}{%
\subsection{Keyword}\label{keyword-1}}

SHEDS

\hypertarget{details-4}{%
\subsection{Details}\label{details-4}}

This function obtains chemical-specific properties from the chem.data
data set.That data set contains point value for each variable, which
this function defines distributions around, with the exception of
molecular weight. The constructed distributions reflect both uncertainty
and variability (for example, vapor pressure will vary with temperature,
humidity, and air pressure or altitude, and may depend on the product
formulation). For each variable, a random sample is generated for each
simulated person, and the set of variables becomes the output data set.
Input surface loading variables are in units of ug/cm2 or ug/cm2/day and
are converted by the function to meters, to avoid the need for
conversion factors in later equations. These conversions are done after
the random sampling, as otherwise the correct conversions depend on the
distributional form (for example, par2 is changed for the normal, but
not for the lognormal).

\hypertarget{author-4}{%
\subsection{Author}\label{author-4}}

Kristin Isaacs, Graham Glen

\hypertarget{seealso-4}{%
\subsection{Seealso}\label{seealso-4}}

\texttt{{[}Fugacity{]}(Fugacity)}, \texttt{{[}Run{]}(Run)},
\texttt{{[}run{]}(run)}, \texttt{{[}Chem\_props{]}(Chem\_props)},
\texttt{{[}get.fug.concs{]}(get.fug.concs)}

\hypertarget{value-4}{%
\subsection{Value}\label{value-4}}

samples A data set with the chemical specific parameters for each
combination of chemical and simulated person. For each chemical, the
chemical specific parameters assigned to a given person are randomly
sampled from distributions on those parameters. These distributions are
created from point estimates to reflect real-world uncertainty and
variability.

\hypertarget{chem.scenarios-chem.scenarios}{%
\section{\texorpdfstring{\texttt{chem.scenarios}:
chem.scenarios}{chem.scenarios: chem.scenarios}}\label{chem.scenarios-chem.scenarios}}

\hypertarget{description-9}{%
\subsection{Description}\label{description-9}}

This function summarizes the all.scenarios data set by condensing each
chemical-scenario combination into one row. The number of rows of data
for this combination on all.scenarios is recorded here.

\hypertarget{usage-9}{%
\subsection{Usage}\label{usage-9}}

\begin{Shaded}
\begin{Highlighting}[]
\FunctionTok{chem.scenarios}\NormalTok{(all)}
\end{Highlighting}
\end{Shaded}

\hypertarget{arguments-6}{%
\subsection{Arguments}\label{arguments-6}}

\begin{itemize}
\tightlist
\item
  \texttt{all}: This is the master list of all chemicals, and all
  exposure scenarios specific to each chemical, to be evaluated in the
  current model run. This data set is compiled internally according to
  the user's specifications on the \texttt{{[}Run\_85687{]}(Run\_85687)}
  file. The user may create multiple scenarios files for special
  purposes, for example, for selected chemical classes, or selected
  exposure pathways. A model run consists of two nested loops: the outer
  loop over chemicals and the inner loop over the scenarios specific to
  that chemical.
\end{itemize}

\hypertarget{author-5}{%
\subsection{Author}\label{author-5}}

Kristin Isaacs, Graham Glen

\hypertarget{seealso-5}{%
\subsection{Seealso}\label{seealso-5}}

\texttt{{[}Run{]}(Run)},
\texttt{{[}Source\_chem\_foods{]}(Source\_chem\_foods)},
\texttt{{[}Source\_chem\_prods{]}(Source\_chem\_prods)},
\texttt{{[}Source\_scen\_food{]}(Source\_scen\_food)},
\texttt{{[}Source\_scen\_prods{]}(Source\_scen\_prods)}

\hypertarget{value-5}{%
\subsection{Value}\label{value-5}}

\hypertarget{combine_output-combine_output}{%
\section{\texorpdfstring{\texttt{combine\_output}:
combine\_output}{combine\_output: combine\_output}}\label{combine_output-combine_output}}

\hypertarget{description-10}{%
\subsection{Description}\label{description-10}}

combine\_output is a post-processing tool that extracts selected
statistics (those specified in the `metrics' argument) from the
`allstats' file for each chemical in a single model run, and combines
them onto out.file.

\hypertarget{usage-10}{%
\subsection{Usage}\label{usage-10}}

\begin{Shaded}
\begin{Highlighting}[]
\FunctionTok{combine\_output}\NormalTok{(}
  \AttributeTok{run.name =}\NormalTok{ specs}\SpecialCharTok{$}\NormalTok{run.name,}
  \AttributeTok{out.file =} \StringTok{"SHEDSOutCombined.csv"}\NormalTok{,}
  \AttributeTok{metrics =} \FunctionTok{c}\NormalTok{(}\StringTok{"5\%"}\NormalTok{, }\StringTok{"50\%"}\NormalTok{, }\StringTok{"75\%"}\NormalTok{, }\StringTok{"95\%"}\NormalTok{, }\StringTok{"99\%"}\NormalTok{, }\StringTok{"mean"}\NormalTok{, }\StringTok{"sd"}\NormalTok{)}
\NormalTok{)}
\end{Highlighting}
\end{Shaded}

\hypertarget{arguments-7}{%
\subsection{Arguments}\label{arguments-7}}

\begin{itemize}
\tightlist
\item
  \texttt{run.name}: default = name of last run (in current R session)
\item
  \texttt{out.file}: default=``SHEDSOutCombined.csv''
\item
  \texttt{metrics}: the summary statsistics to be pulled. default =
  c(``5\%'', ``50\%'', ``75\%'', ``95\%'', ``99\%'',``mean'',``sd'')
\end{itemize}

\hypertarget{details-5}{%
\subsection{Details}\label{details-5}}

In a multichemical run, SHEDS creates separate output files for each
chemical. If the analyst wants to compare exposure or dose metrics
across chemicals, use combine\_output to put the relevant information
for all chemicals together on one file.

\hypertarget{author-6}{%
\subsection{Author}\label{author-6}}

Kristin Isaacs, Graham Glen

\hypertarget{value-6}{%
\subsection{Value}\label{value-6}}

finaldata A R object that has ``exp.dermal'', ``exp.ingest'',
m''exp.inhal'', ``dose.inhal'', ``dose.intake'', ``abs.dermal.ug'',
``abs.ingest.ug'', ``abs.inhal.ug'', ``abs.tot.ug'', ``abs.tot.mgkg'',
``ddd.mass'' for the chemcial of question. This is tablated for the
``5\%'', ``50\%'', ``75\%'', ``95\%'', ``99\%'', quantile of exposure as
well as the mean and sd
``finaldata,paste0(''output/``,run.name,''/``,out.file.csv'') this is
the CSV of finaldata

\hypertarget{create.scen.factors-create.scen.factors}{%
\section{\texorpdfstring{\texttt{create.scen.factors}:
create.scen.factors}{create.scen.factors: create.scen.factors}}\label{create.scen.factors-create.scen.factors}}

\hypertarget{description-11}{%
\subsection{Description}\label{description-11}}

This function takes the information on distributions from the
\texttt{all.scenarios} data set (which comes from the
\texttt{{[}source\_variables\_12112015{]}(source\_variables\_12112015)}
input file) and converts it into the parameter set needed by SHEDS.HT.

\hypertarget{usage-11}{%
\subsection{Usage}\label{usage-11}}

\begin{Shaded}
\begin{Highlighting}[]
\FunctionTok{create.scen.factors}\NormalTok{(f)}
\end{Highlighting}
\end{Shaded}

\hypertarget{arguments-8}{%
\subsection{Arguments}\label{arguments-8}}

\begin{itemize}
\tightlist
\item
  \texttt{f}: An internally generated data set from the
  \texttt{{[}Source\_vars{]}(Source\_vars)} input file (as specified in
  the \texttt{{[}Run{]}(Run)} file) containing data on the distribution
  of each source variable in SHEDS.HT.
\end{itemize}

\hypertarget{details-6}{%
\subsection{Details}\label{details-6}}

The steps involved in this function are 1) converting
\texttt{prevalence} from a percentage to a binomial form, 2) converting
\texttt{CV} to standard deviation for normals, and 3) converting
\texttt{Mean} and \texttt{CV} to par1 and par2 for lognormals. Note that
\texttt{prevalence} in SHEDS becomes a binomial distribution, which
returns a value of either 0 or 1 when evaluated. Each simulated person
either ``does'' or ``does not'' partake in this scenario. Similar logic
applies to the \texttt{frequency}variable, except that the returned
values may be larger than one (that is, 2 or more) for very frequent
scenarios. All the exposure equations contain the \texttt{prevalence}
variable. If \texttt{prevalence} is set to one for that person, the
exposure is as expected, but if \texttt{prevalence}=0 then no exposure
occurs.

\hypertarget{author-7}{%
\subsection{Author}\label{author-7}}

Kristin Isaacs, Graham Glen

\hypertarget{seealso-6}{%
\subsection{Seealso}\label{seealso-6}}

\texttt{{[}Source\_vars{]}(Source\_vars)}, \texttt{{[}Run{]}(Run)}

\hypertarget{value-7}{%
\subsection{Value}\label{value-7}}

dt A data set with values and probabilities associated with each
variable distribution presented in \texttt{f}. All variables in
\texttt{f} are also retained.

\hypertarget{diet.diary.pools-diet.diary.pools}{%
\section{\texorpdfstring{\texttt{diet.diary.pools}:
diet.diary.pools}{diet.diary.pools: diet.diary.pools}}\label{diet.diary.pools-diet.diary.pools}}

\hypertarget{description-12}{%
\subsection{Description}\label{description-12}}

Assigns activity diaries from the
\texttt{{[}Diet\_diaries{]}(Diet\_diaries)} input on the
\texttt{{[}Run{]}(Run)} file (read through the
\texttt{{[}read.run.file{]}(read.run.file)} function) to pools based on
age and gender.

\hypertarget{usage-12}{%
\subsection{Usage}\label{usage-12}}

\begin{Shaded}
\begin{Highlighting}[]
\FunctionTok{diet.diary.pools}\NormalTok{(diaries, specs)}
\end{Highlighting}
\end{Shaded}

\hypertarget{arguments-9}{%
\subsection{Arguments}\label{arguments-9}}

\begin{itemize}
\tightlist
\item
  \texttt{specs}: Output of the
  \texttt{{[}read.run.file{]}(read.run.file)} function, which can be
  modified by the \texttt{{[}update.specs{]}(update.specs)} function
  before input into \texttt{diet.diary.pools}.
\item
  \texttt{diet.diaries}: A data set created internally in SHEDS.HT
  through the \texttt{{[}read.diet.diaries{]}(read.diet.diaries)}
  function. The data are daily diaries of dietary consumption by food
  group. Each line represents one person-day, with demographic variables
  followed by amounts (in grams/day) for a list of food types indicated
  by a short abbreviation on the header line.
\end{itemize}

\hypertarget{details-7}{%
\subsection{Details}\label{details-7}}

The \texttt{diet.diary.pools} function assigns the dietary diaries to
pools. A given diary may belong to many pools, since every year of age
has its own pool. Large pools may contain a list of several hundred
diary numbers. The code contains four loops and on each step performs a
sub-setting of the list of diary numbers.

\hypertarget{note-1}{%
\subsection{Note}\label{note-1}}

The input to \texttt{diet.diary.pools}, \texttt{diet.diaries}, is
created by reading in the
\texttt{{[}Diet\_diaries{]}(Diet\_diaries)}input file (specified on the
\texttt{{[}Run{]}(Run)} file) with the
\texttt{{[}read.diet.diaries{]}(read.diet.diaries)} function within the
\texttt{{[}run{]}(run)} function.

\hypertarget{author-8}{%
\subsection{Author}\label{author-8}}

Kristin Isaacs, Graham Glen

\hypertarget{seealso-7}{%
\subsection{Seealso}\label{seealso-7}}

\texttt{{[}Diet\_diaries{]}(Diet\_diaries)}, \texttt{{[}run{]}(run)},
\texttt{{[}read.run.file{]}(read.run.file)}, \texttt{{[}Run{]}(Run)}

\hypertarget{value-8}{%
\subsection{Value}\label{value-8}}

dpool A vector of lists. The length of \texttt{dpool} is the product of
the gender and age inputs specified in the \texttt{{[}Run{]}(Run)} file.
Each element is a list of acceptable diet diary numbers for each year of
age and each gender combination. In addition, if the function runs
successfully, the following message will be printed: ``Dietary Diary
Pooling completed''

\hypertarget{dir.dermal-dir.dermal}{%
\section{\texorpdfstring{\texttt{dir.dermal}:
dir.dermal}{dir.dermal: dir.dermal}}\label{dir.dermal-dir.dermal}}

\hypertarget{description-13}{%
\subsection{Description}\label{description-13}}

Models the dermal exposure scenario for each theoretical person.

\hypertarget{usage-13}{%
\subsection{Usage}\label{usage-13}}

\begin{Shaded}
\begin{Highlighting}[]
\FunctionTok{dir.dermal}\NormalTok{(sd, cd)}
\end{Highlighting}
\end{Shaded}

\hypertarget{arguments-10}{%
\subsection{Arguments}\label{arguments-10}}

\begin{itemize}
\tightlist
\item
  \texttt{sd}: The chemical-scenario data specific to relevant
  combinations of chemical and scenario. Generated internally.
\item
  \texttt{cd}: The list of scenario-specific information for the
  chemicals being evaluated. Generated internally.
\end{itemize}

\hypertarget{details-8}{%
\subsection{Details}\label{details-8}}

The dermal exposure scenario is relatively straightforward. The function
produces a prevalence value, which reflects the fraction of the
population who use this scenario at all. It also produces a frequency
value, which is the mean number of times per year this scenario occurs
among that fraction of the population specified by prevalence. Since
SHEDS operates on the basis of one random day, the frequency is divided
by 365 and then passed to the \texttt{{[}p.round{]}(p.round)}
(probabilistic rounding) function, which rounds either up or down to the
nearest integer. Very common events may happen more than once in a day.
The function also produces a mass variable, which refers to the mass of
the product in grams in a typical usage event. The composition is the
percentage of that mass that is the chemical in question. The resid
variable measures the fraction that is likely to remain on the skin when
the usage event ends. The final output is the total dermal exposure for
each chemical-individual combination in micrograms, which is the product
of the above variables multiplied by a factor of 1E+06.

\hypertarget{author-9}{%
\subsection{Author}\label{author-9}}

Kristin Isaacs, Graham Glen

\hypertarget{seealso-8}{%
\subsection{Seealso}\label{seealso-8}}

\texttt{{[}run{]}(run)}, \texttt{{[}p.round{]}(p.round)}

\hypertarget{value-9}{%
\subsection{Value}\label{value-9}}

dir.derm For each person, the calculated direct dermal exposure that
occurs when a chemical-containing product is used. This does not include
later contact with treated objects, which is indirect exposure.

\hypertarget{dir.ingested-dir.ingested}{%
\section{\texorpdfstring{\texttt{dir.ingested}:
dir.ingested}{dir.ingested: dir.ingested}}\label{dir.ingested-dir.ingested}}

\hypertarget{description-14}{%
\subsection{Description}\label{description-14}}

Models the ingestion exposure scenario for each theoretical person.

\hypertarget{usage-14}{%
\subsection{Usage}\label{usage-14}}

\begin{Shaded}
\begin{Highlighting}[]
\FunctionTok{dir.ingested}\NormalTok{(sd, cd)}
\end{Highlighting}
\end{Shaded}

\hypertarget{arguments-11}{%
\subsection{Arguments}\label{arguments-11}}

\begin{itemize}
\tightlist
\item
  \texttt{sd}: The chemical-scenario data specific to relevant
  combinations of chemical and scenario. Generated internally.
\item
  \texttt{cd}: The list of scenario-specific information for the
  chemicals being evaluated. Generated internally.
\end{itemize}

\hypertarget{details-9}{%
\subsection{Details}\label{details-9}}

This scenario is for accidental ingestion during product usage. Typical
examples are toothpaste, mouthwash, lipstick or chap stick, and similar
products used on the face or mouth. The function produces a
\texttt{prevalence} value, which reflects the fraction of the population
who use this scenario at all. It also produces a \texttt{frequency}
value, which is the mean number of times per year this scenario occurs
among that fraction of the population specified by prevalence. Since
SHEDS operates on the basis of one random day, the frequency is divided
by 365 and then passed to the \texttt{{[}p.round{]}(p.round)}
(probabilistic rounding) function, which rounds either up or down to the
nearest integer. Very common events may happen more than once in a day.
The function also produces a \texttt{mass} variable, which refers to the
mass of the product in grams in a typical usage event. The
\texttt{composition} is the percentage of that mass that is the chemical
in question. The \texttt{ingested} variable represents the percentage of
the mass applied that becomes ingested. Since these products are not
intended to be swallowed, this should typically be quite small (under
5\%). The final output is the total incidental ingested exposure for
each chemical-individual combination in micrograms, which is the product
of the above variables multiplied by a factor of 1E6.

\hypertarget{author-10}{%
\subsection{Author}\label{author-10}}

Kristin Isaacs, Graham Glen

\hypertarget{value-10}{%
\subsection{Value}\label{value-10}}

dir.ingest For each person, the calculated quantity of a given chemical
incidentally ingested during or immediately after use of products such
as toothpaste. Does not include exposure via food and drinking water.

\hypertarget{dir.inhal.aer-dir.inhal.aer}{%
\section{\texorpdfstring{\texttt{dir.inhal.aer}:
dir.inhal.aer}{dir.inhal.aer: dir.inhal.aer}}\label{dir.inhal.aer-dir.inhal.aer}}

\hypertarget{description-15}{%
\subsection{Description}\label{description-15}}

Models the inhalation exposure from the use of aerosol products for each
theoretical person.

\hypertarget{usage-15}{%
\subsection{Usage}\label{usage-15}}

\begin{Shaded}
\begin{Highlighting}[]
\FunctionTok{dir.inhal.aer}\NormalTok{(sd, cd, cb, io)}
\end{Highlighting}
\end{Shaded}

\hypertarget{arguments-12}{%
\subsection{Arguments}\label{arguments-12}}

\begin{itemize}
\tightlist
\item
  \texttt{sd}: The chemical-scenario data specific to relevant
  combinations of chemical and scenario. Generated internally.
\item
  \texttt{cd}: The list of scenario-specific information for the
  chemicals being evaluated. Generated internally.
\item
  \texttt{cb}: A copy of the \texttt{base} data set output from the
  \texttt{{[}make.cbase{]}(make.cbase)} function, with columns added for
  exposure variables.
\item
  \texttt{io}: A binary variable indicating whether the volume of the
  aerosol is used to approximate the affected volume.
\end{itemize}

\hypertarget{details-10}{%
\subsection{Details}\label{details-10}}

This scenario considers inhalation exposure from the use of aerosol
products. Typical examples include hairspray and spray-on mosquito
repellent. The function produces a \texttt{prevalence} value, which
reflects the fraction of the population who use this scenario at all. It
also produces a \texttt{frequency} value, which is the mean number of
times per year this scenario occurs among that fraction of the
population specified by prevalence. Since SHEDS operates on the basis of
one random day, the \texttt{frequency} is divided by 365 and then passed
to the \texttt{{[}p.round{]}(p.round)} (probabilistic rounding)
function, which rounds either up or down to the nearest integer. Very
common events may happen more than once in a day. The function also
produces a \texttt{mass} variable, which refers to the mass of the
product in grams in a typical usage event. The \texttt{composition} is
the percentage of that mass that is the chemical in question.
\texttt{frac.aer} is the fraction of the product mass that becomes
aerosolized, and the \texttt{volume} affected by the use is approximated
to allow the calculation of a concentration or density. Defaults are set
in the code if these variables are missing from the input file. Exposure
for the inhalation pathway has units of micrograms per cubic meter,
reflecting the average air concentration of the chemical. An
\texttt{airconc} variable is defined using \texttt{mass},
\texttt{composition}, \texttt{frac.aer}, and \texttt{volume}. Since
exposure depends on the time-averaged concentration, a duration is
necessary. For example, if one spends five minutes in an aerosol cloud
and the rest of the day in clean air, the daily exposure is the cloud
concentration multiplied by 5/1440 (where 1440 is the number of minutes
in a day). This function also calculates the inhaled dose, in units of
micrograms per day. The dose equals the product of \texttt{exposure}
(g/m3), basal ventilation rate, \texttt{bvr} (m3/day), the METS factor
of 1.75 (typically people inhale air at an average of 1.75 times the
basal rate to support common daily activities), and a conversion factor
of 1E+06 from grams to micrograms.

\hypertarget{author-11}{%
\subsection{Author}\label{author-11}}

Kristin Isaacs, Graham Glen keyword\textasciitilde SHEDS

\hypertarget{value-11}{%
\subsection{Value}\label{value-11}}

dir.inh.aer The calculated quantity of chemical inhalation from
aerosols, like hairspray and similar products, that are directly
injected into the air on or around each exposed person.

\hypertarget{dir.inhal.vap-dir.inhal.vap}{%
\section{\texorpdfstring{\texttt{dir.inhal.vap}:
dir.inhal.vap}{dir.inhal.vap: dir.inhal.vap}}\label{dir.inhal.vap-dir.inhal.vap}}

\hypertarget{description-16}{%
\subsection{Description}\label{description-16}}

Models the inhalation exposure from the vapors of volatile chemicals for
each theoretical person.

\hypertarget{usage-16}{%
\subsection{Usage}\label{usage-16}}

\begin{Shaded}
\begin{Highlighting}[]
\FunctionTok{dir.inhal.vap}\NormalTok{(sd, cd, cprops, cb, io)}
\end{Highlighting}
\end{Shaded}

\hypertarget{arguments-13}{%
\subsection{Arguments}\label{arguments-13}}

\begin{itemize}
\tightlist
\item
  \texttt{sd}: The chemical-scenario data specific to relevant
  combinations of chemical and scenario. Generated internally.
\item
  \texttt{cd}: The list of scenario-specific information for the
  chemicals being evaluated. Generated internally.
\item
  \texttt{cprops}: The chemical properties required for SHEDS-HT. The
  default file (the
  \texttt{{[}Chem\_props\ file{]}(Chem\_props\%20file)} read in by the
  \texttt{{[}read.chem.props{]}(read.chem.props)} function and modified
  before input into the current function) was prepared from publicly
  available databases using a custom program (not part of SHEDS-HT). The
  default file contains 7 numerical inputs per chemical, and the
  required properties are molecular weight (\texttt{MW}), vapor pressure
  (\texttt{VP.Pa}), solubility (\texttt{water.sol.mg.l}), octanol-water
  partition coefficient (\texttt{log.Kow}), air decay rate
  (\texttt{half.air.hr}), decay rate on surfaces
  (\texttt{half.sediment.hr}), and permeability coefficient
  (\texttt{Kp}).
\item
  \texttt{cb}: A copy of the \texttt{base} data set output from the
  \texttt{{[}make.cbase{]}(make.cbase)} function, with columns added for
  exposure variables.
\item
  \texttt{io}: A binary variable indicating whether the volume of the
  aerosol is used to approximate the affected volume.
\end{itemize}

\hypertarget{details-11}{%
\subsection{Details}\label{details-11}}

This scenario considers inhalation exposure from vapors (not aerosols).
For example, painting will result in the inhalation of vapor, but it
does not involve aerosols (unless it is spray paint). For this scenario,
the vapor pressure and the molecular weight are relevant variables for
determining exposure. These variables are included in the input to the
\texttt{cprops} argument, which is drawn internally from the
\texttt{{[}Chem\_props{]}(Chem\_props)} file. The function produces a
\texttt{prevalence} value, which reflects the fraction of the population
who use this scenario at all. It also produces a \texttt{frequency}
value, which is the mean number of times per year this scenario occurs
among that fraction of the population specified by \texttt{prevalence}.
Since SHEDS operates on the basis of one random day, the
\texttt{frequency} is divided by 365 and then passed to the
\texttt{{[}p.round{]}(p.round)} (probabilistic rounding) function, which
rounds either up or down to the nearest integer. Very common events may
happen more than once in a day. The function also produces a
\texttt{mass} variable, which refers to the mass of the product in grams
in a typical usage event. The \texttt{composition} is the percentage of
that mass that is the chemical in question. The \texttt{evap} variable
is an effective evaporated mass, calculated using the \texttt{mass},
\texttt{composition} (converted from percent to a fraction), the vapor
pressure as a surrogate for partial pressure, and \texttt{duration} of
product use. The \texttt{duration} term is made unitless by dividing by
5 (minutes), which is an assumed time constant. The effective air
concentation \texttt{airconc} is calculated as
\texttt{evap}/\texttt{volume}. The value for \texttt{airconce}is capped
by \texttt{maxconc}, which represents the point at which evaporation
ceases. For chemicals used for a short duration, or with low vapor
presssure, \texttt{maxconc} might not be reached before usage stops.
Once \texttt{airconc} is established, the function also calculates the
inhaled dose, in units of micrograms per day. The dose equals the
product of \texttt{exposure} (g/m3), basal ventilation rate,
\texttt{bvr} (m3/day), the METS factor of 1.75 (typically people inhale
air at an average of 1.75 times the basal rate to support common daily
activities), and a conversion factor of 1E6 from grams to micrograms.

\hypertarget{author-12}{%
\subsection{Author}\label{author-12}}

Kristin Isaacs, Graham Glen

\hypertarget{seealso-9}{%
\subsection{Seealso}\label{seealso-9}}

\texttt{{[}run{]}(run)}, \texttt{{[}p.round{]}(p.round)},
\texttt{{[}read.chem.props{]}(read.chem.props)},
\texttt{{[}Chem\_props{]}(Chem\_props)}

\hypertarget{distrib-distrib}{%
\section{\texorpdfstring{\texttt{distrib}:
distrib}{distrib: distrib}}\label{distrib-distrib}}

\hypertarget{description-17}{%
\subsection{Description}\label{description-17}}

produces random samples from distributuions

\hypertarget{usage-17}{%
\subsection{Usage}\label{usage-17}}

\begin{Shaded}
\begin{Highlighting}[]
\FunctionTok{distrib}\NormalTok{(}
  \AttributeTok{shape =} \StringTok{""}\NormalTok{,}
  \AttributeTok{par1 =} \ConstantTok{NA}\NormalTok{,}
  \AttributeTok{par2 =} \ConstantTok{NA}\NormalTok{,}
  \AttributeTok{par3 =} \ConstantTok{NA}\NormalTok{,}
  \AttributeTok{par4 =} \ConstantTok{NA}\NormalTok{,}
  \AttributeTok{lt =} \ConstantTok{NA}\NormalTok{,}
  \AttributeTok{ut =} \ConstantTok{NA}\NormalTok{,}
  \AttributeTok{resamp =} \StringTok{"y"}\NormalTok{,}
  \AttributeTok{n =} \DecValTok{1}\NormalTok{,}
  \AttributeTok{q =} \ConstantTok{NA}\NormalTok{,}
  \AttributeTok{p =} \FunctionTok{c}\NormalTok{(}\DecValTok{1}\NormalTok{),}
  \AttributeTok{v =} \StringTok{""}
\NormalTok{)}
\end{Highlighting}
\end{Shaded}

\hypertarget{arguments-14}{%
\subsection{Arguments}\label{arguments-14}}

\begin{itemize}
\tightlist
\item
  \texttt{shape}: Required with no default. The permited inputs are
  Bernoulli, binomial, beta, discrete, empirical, exponential, gamma,
  lognormal, normal, point, probability, triangle, uniform, and Weibull.
\item
  \texttt{par1}: optional value requared for some shapes. Defulat = none
\item
  \texttt{par2}: optional value requared for some shapes. Defulat = none
\item
  \texttt{par3}: optional value requared for some shapes. Defulat = none
\item
  \texttt{par4}: optional value requared for some shapes. Defulat = none
\item
  \texttt{resamp}: optional value requared for some shapes. Defulat =
  `y'
\item
  \texttt{n}: Optional Default = 1
\item
  \texttt{q}: Optional Default = none
\item
  \texttt{p}: Optional Default = c(1)
\item
  \texttt{v}: Optional Default = none
\item
  \texttt{It}: optional value requared for some shapes. Defulat = none
\item
  \texttt{Ut}: optional value requared for some shapes. Defulat = none
\end{itemize}

\hypertarget{details-12}{%
\subsection{Details}\label{details-12}}

This process is central to SHEDS because it is a stochastic model. The
``shape'' is the essential argument, with others being required for
certain shapes. For all shapes, one of ``n'' or ``q'' must be specified.
If ``n'' is given, then Distrib returns a vector of ``n'' independent
random samples from the specified distribution. If ``q'' is given, it
must be a vector of numeric values, each between zero and one. These are
interpreted as the quantiles of the distribution to be returned. When
``q'' is given, Distrib does not generate any random values, it just
evaluates the requested quantiles. The empirical shape requires argument
``v'' as a list of possible values to be returned, each with equal
probability. The other shapes require one or more of par1-par4 to be
specified. See the SHEDS Technical Manual for more details on the
meanings of par1-par4, which vary by shape. ``Lower.trun'' is the lower
truncation point, meaning the smallest value that can be returned.
Similarly, ``upper.trun'' is the largest value that may be returned. Not
all distributions use lower.trun and/or upper.trun, but they should be
specified for unbounded shapes like the Normal distribution. ``Resamp''
is a flag to indicate the resolution for generated values outside the
truncation limits. If resamp=``yes'' then effectively new values are
generated until they are within the limits. If resamp=``no'', values
outside the limits are moved to those limits. ``P'' is a list of
probabilities that are used only with the ``discrete'' or
``probability'' shapes. The ``p'' values are essentially weights for a
list of discrete values that may be returned. The ``empirical''
distribution also returns discrete values, but assigns them equal
weights, so then ``p'' is not needed.

\hypertarget{author-13}{%
\subsection{Author}\label{author-13}}

Kristin Isaacs, Graham Glen

\hypertarget{value-12}{%
\subsection{Value}\label{value-12}}

A vector of ``n'' values from one distribution, where ``n'' is either
the input argument (if given), or the length of the input vector ``q''.

\hypertarget{down.the.drain.mass-down.the.drain.mass}{%
\section{\texorpdfstring{\texttt{down.the.drain.mass}:
down.the.drain.mass}{down.the.drain.mass: down.the.drain.mass}}\label{down.the.drain.mass-down.the.drain.mass}}

\hypertarget{description-18}{%
\subsection{Description}\label{description-18}}

Models the quantity of chemical entering the waste water system on a per
person-day basis.

\hypertarget{usage-18}{%
\subsection{Usage}\label{usage-18}}

\begin{Shaded}
\begin{Highlighting}[]
\FunctionTok{down.the.drain.mass}\NormalTok{(sd, cd)}
\end{Highlighting}
\end{Shaded}

\hypertarget{arguments-15}{%
\subsection{Arguments}\label{arguments-15}}

\begin{itemize}
\tightlist
\item
  \texttt{sd}: The chemical-scenario data specific to relevant
  combinations of chemical and scenario. Generated internally.
\item
  \texttt{cd}: The list of scenario-specific information for the
  chemicals being evaluated. Generated internally.
\end{itemize}

\hypertarget{details-13}{%
\subsection{Details}\label{details-13}}

This function models the simplest of all the current scenarios. It
evaluates the amount of chemical entering the waste water system, on a
per person-day basis. The ``exposure'' is to a system, not a person, but
this method uses one person's actions to estimate their contribution to
the total. The function produces a \texttt{prevalence} value, which
reflects the fraction of the population who use this scenario at all. It
also produces a \texttt{frequency} value, which is the mean number of
times per year this scenario occurs among that fraction of the
population specified by prevalence. Since SHEDS operates on the basis of
one random day, the \texttt{frequency} is divided by 365 and then passed
to the \texttt{{[}p.round{]}(p.round)} (probabilistic rounding)
function, which rounds either up or down to the nearest integer. Very
common events may happen more than once in a day. The function also
produces a \texttt{mass} variable, which refers to the mass of the
product in grams in a typical usage event. The \texttt{composition} is
the percentage of that mass that is the chemical in question. The final
output, \texttt{exp.ddd.mass}, is the product of the
\texttt{prevalence}, \texttt{frequency}, \texttt{mass},
\texttt{composition}, and the fraction going down the drain
(\texttt{f.drain}, a variable in the \texttt{sd} input). The result is
in grams per person-day.

\hypertarget{author-14}{%
\subsection{Author}\label{author-14}}

Kristin Isaacs, Graham Glen

\hypertarget{seealso-10}{%
\subsection{Seealso}\label{seealso-10}}

\texttt{{[}run{]}(run)}, \texttt{{[}p.round{]}(p.round)}

\hypertarget{value-13}{%
\subsection{Value}\label{value-13}}

exp.ddd.mass The calculated quantity of chemical going down the drain
(i.e., from laundry detergent) and entering the sewer system per person
per day.

\hypertarget{eval.factors-eval.factors}{%
\section{\texorpdfstring{\texttt{eval.factors}:
eval.factors}{eval.factors: eval.factors}}\label{eval.factors-eval.factors}}

\hypertarget{description-19}{%
\subsection{Description}\label{description-19}}

Assigns distributions for each relevant exposure factor for each
theoretical persons to be modeled in SHEDS.HT.

\hypertarget{usage-19}{%
\subsection{Usage}\label{usage-19}}

\begin{Shaded}
\begin{Highlighting}[]
\FunctionTok{eval.factors}\NormalTok{(r, q, ef)}
\end{Highlighting}
\end{Shaded}

\hypertarget{arguments-16}{%
\subsection{Arguments}\label{arguments-16}}

\begin{itemize}
\tightlist
\item
  \texttt{r}: Vector of row numbers, equivalent to the number of
  theoretical people in the model sample. This input is created
  internally by the \texttt{{[}add.factors{]}(add.factors)} function.
\item
  \texttt{q}: User-specified list of desired quantiles to be included in
  the output.
\item
  \texttt{ef}: Distributional parameters for the exposure factors; an
  output of the \texttt{{[}exp.factors{]}(exp.factors)} function and an
  input argument to the \texttt{{[}add.factors{]}(add.factors)} function
\end{itemize}

\hypertarget{author-15}{%
\subsection{Author}\label{author-15}}

Krisin Isaacs, Graham Glen

\hypertarget{seealso-11}{%
\subsection{Seealso}\label{seealso-11}}

\texttt{{[}add.factors{]}(add.factors)},
\texttt{{[}exp.factors{]}(exp.factors)}

\hypertarget{value-14}{%
\subsection{Value}\label{value-14}}

z A data frame specifying the form of the distribution and relevant
parameters for each combination of exposure factor and individual
person. The parameters include:

\begin{itemize}
\tightlist
\item
  \texttt{form}: The form of the distribution for each exposure factor
  (i.e., point, triangle).
\item
  \texttt{par1-par4}: The parameters associated with the distributional
  form specified in form.
\item
  \texttt{lower.trun}: the lower limit of values to be included in the
  distribution.
\item
  \texttt{upper.trun}: The upper limit of values to be included in the
  distribution.
\item
  \texttt{resamp}: Logical field where: yes=resample, no=stack at
  truncation bounds.
\item
  \texttt{q}: The quantiles of the distribution corresponding to those
  specified by the user in the \texttt{q} input.
\item
  \texttt{p}: The probabilities associated with the quantiles.
\item
  \texttt{v}: The values associated with the quantiles.
\end{itemize}

\hypertarget{exportdatatables-a-function-used-to-export-.rda-files-distributed-with-shedsht-package}{%
\section{\texorpdfstring{\texttt{ExportDataTables}: A function used to
export .rda files distributed with ShedsHT
package}{ExportDataTables: A function used to export .rda files distributed with ShedsHT package}}\label{exportdatatables-a-function-used-to-export-.rda-files-distributed-with-shedsht-package}}

\hypertarget{description-20}{%
\subsection{Description}\label{description-20}}

This is a utility function used to export consumer product data sets
into CSV format.

\hypertarget{usage-20}{%
\subsection{Usage}\label{usage-20}}

\begin{Shaded}
\begin{Highlighting}[]
\FunctionTok{ExportDataTables}\NormalTok{(}
\NormalTok{  data\_set\_name,}
\NormalTok{  output\_pth,}
\NormalTok{  output\_fname,}
  \AttributeTok{quote\_opt =} \ConstantTok{TRUE}\NormalTok{,}
  \AttributeTok{na\_opt =} \StringTok{""}\NormalTok{,}
  \AttributeTok{row\_names\_opt =} \ConstantTok{FALSE}
\NormalTok{)}
\end{Highlighting}
\end{Shaded}

\hypertarget{arguments-17}{%
\subsection{Arguments}\label{arguments-17}}

\begin{itemize}
\tightlist
\item
  \texttt{data\_set\_name}: A string type parameter, represents name of
  an ``.rda'' data set.
\item
  \texttt{output\_pth}: A string type parameter, represents location to
  store exported CSV file.
\item
  \texttt{output\_fname}: A string type parameter, represents name of
  the exported CSV file.
\item
  \texttt{quote\_opt}: A string type parameter, represents a logical
  argument for if any character or factor columns should be surrounded
  by double quotes or not.
\item
  \texttt{na\_opt}: A string type parameter, the string to use for
  missing values in the data.
\item
  \texttt{row\_names\_opt}: A string type parameter, either a logical
  value indicating whether the row names of x are to be written along
  with x, or a character vector of row names to be written.
\end{itemize}

\hypertarget{value-15}{%
\subsection{Value}\label{value-15}}

No variable will be returned. Instead, the function will save the data
object into the file of choice.

\hypertarget{filter_sources-filter_sources}{%
\section{\texorpdfstring{\texttt{filter\_sources}:
filter\_sources}{filter\_sources: filter\_sources}}\label{filter_sources-filter_sources}}

\hypertarget{description-21}{%
\subsection{Description}\label{description-21}}

This function is used to select a subset of sources from another
``source'' file.

\hypertarget{usage-21}{%
\subsection{Usage}\label{usage-21}}

\begin{Shaded}
\begin{Highlighting}[]
\FunctionTok{filter\_sources}\NormalTok{(}
  \AttributeTok{in.file =} \StringTok{"source\_scen\_prods.csv"}\NormalTok{,}
  \AttributeTok{out.file =} \StringTok{"source\_scen\_subset.csv"}\NormalTok{,}
  \AttributeTok{ids =} \FunctionTok{c}\NormalTok{(}\StringTok{"ALL"}\NormalTok{),}
  \AttributeTok{types =} \FunctionTok{c}\NormalTok{(}\StringTok{"ALL"}\NormalTok{)}
\NormalTok{)}
\end{Highlighting}
\end{Shaded}

\hypertarget{arguments-18}{%
\subsection{Arguments}\label{arguments-18}}

\begin{itemize}
\tightlist
\item
  \texttt{In.file}: a csv file that contains expsoure info default =
  ``source\_scen\_prods.csv''
\item
  \texttt{Out.file}: default = ``source\_scen\_prods.csv''
\item
  \texttt{IDS}: default = ``ALL''
\item
  \texttt{Types}: default = ``ALL''
\end{itemize}

\hypertarget{details-14}{%
\subsection{Details}\label{details-14}}

This function is used to select a subset of sources from another
``source'' file. This is achieved by specifying either a list of the
desired source ids, or one or more sources types. The allowed types are
``A'' for articles, ``F'' for foods, or ``P'' for products. Use the c()
function to list more than one item (e.g.~types=c(``F'',``P'') for foods
and products). Any of the three types of source files (that is,
source\_scen, source\_chem, or source\_vars) may be used. Specifying
specific sources is achieved using the ``ids'' argument. This may
require examining the in.file beforehand, to obtain the correct
source.id values for the desired sources.

\hypertarget{value-16}{%
\subsection{Value}\label{value-16}}

A .csv file of the same type as in.file, with the same variables and
data, but fewer rows. The selected rows have source.type matching one of
the elements in the ``types'' argument, and also have source.id matching
one of the elements of the ``ids'' argument. If either argument is
missing, then all sources automatically match it. Filter\_sources also
returns an R object containing the same data as the output .csv file.

\hypertarget{food.migration-food.migration}{%
\section{\texorpdfstring{\texttt{food.migration}:
food.migration}{food.migration: food.migration}}\label{food.migration-food.migration}}

\hypertarget{description-22}{%
\subsection{Description}\label{description-22}}

Calculates chemical exposure from migration of chemicals from packaging
or other contact materials into food (which is then consumed). This is
added to the \texttt{dietary} output calculated by the
\texttt{{[}food.residue{]}(food.residue)} function to establish a total
exposure from the dietary pathway.

\hypertarget{usage-22}{%
\subsection{Usage}\label{usage-22}}

\begin{Shaded}
\begin{Highlighting}[]
\FunctionTok{food.migration}\NormalTok{(cdata, sdata, cb, ftype)}
\end{Highlighting}
\end{Shaded}

\hypertarget{arguments-19}{%
\subsection{Arguments}\label{arguments-19}}

\begin{itemize}
\tightlist
\item
  \texttt{cdata}: The list of scenario-specific information for the
  chemicals being evaluated. Generated internally.
\item
  \texttt{sdata}: The chemical-scenario data specific to relevant
  combinations of chemical and scenario. Generated internally.
\item
  \texttt{cb}: A copy of the \texttt{base} data set output from the
  \texttt{{[}make.cbase{]}(make.cbase)} function, with columns added for
  exposure variables.
\item
  \texttt{ftype}: Food consumption database which stores data on
  consumption in grams per day of each food type for each person being
  modeled. Generated internally from the
  \texttt{{[}Diet\_diaries{]}(Diet\_diaries)} data set.
\end{itemize}

\hypertarget{details-15}{%
\subsection{Details}\label{details-15}}

If migration data is stored in the \texttt{cdata} argument, this is
added to the total exposure calculated in the
\texttt{{[}food.residue{]}(food.residue)} function.

\hypertarget{author-16}{%
\subsection{Author}\label{author-16}}

Kristin Isaacs, Graham Glen

\hypertarget{seealso-12}{%
\subsection{Seealso}\label{seealso-12}}

\texttt{{[}Diet\_diaries{]}(Diet\_diaries)}, \texttt{{[}run{]}(run)},
\texttt{{[}p.round{]}(p.round)}, \texttt{{[}make.cbase{]}(make.cbase)},
\texttt{{[}food.residue{]}(food.residue)}

\hypertarget{value-17}{%
\subsection{Value}\label{value-17}}

dietary The calculated quantity of chemical exposure from food residue
and migration of chemicals into food from packaging and other contact
materials in grams per person per day.

\hypertarget{food.residue-food.residue}{%
\section{\texorpdfstring{\texttt{food.residue}:
food.residue}{food.residue: food.residue}}\label{food.residue-food.residue}}

\hypertarget{description-23}{%
\subsection{Description}\label{description-23}}

Models exposure to chemicals from consumption of food containing a known
chemical residue for each theoretical person.

\hypertarget{usage-23}{%
\subsection{Usage}\label{usage-23}}

\begin{Shaded}
\begin{Highlighting}[]
\FunctionTok{food.residue}\NormalTok{(cdata, cb, ftype)}
\end{Highlighting}
\end{Shaded}

\hypertarget{arguments-20}{%
\subsection{Arguments}\label{arguments-20}}

\begin{itemize}
\tightlist
\item
  \texttt{cdata}: The list of scenario-specific information for the
  chemicals being evaluated. Generated internally.
\item
  \texttt{cb}: Output of the \texttt{{[}make.cbase{]}(make.cbase)}
  function.
\item
  \texttt{ftype}: Food consumption database which stores data on
  consumption in grams per day of each food type for each person being
  modeled. Generated internally from the
  \texttt{{[}Diet\_diaries{]}(Diet\_diaries)} data set.
\end{itemize}

\hypertarget{details-16}{%
\subsection{Details}\label{details-16}}

In this function, the variable \texttt{foods} is defined as a list of
the names of the food groups, which are stored in both the
\texttt{ftype} and \texttt{cb} arguments. The FOR loop picks the name of
each food group as a string, and converts it to a variable name. For
example, the food group ``FV'' may have corresponding variables
\texttt{residue.FV}, \texttt{zeros.FV} (number of nondetects), and
\texttt{nonzeros.FV} (number of detects) on the \texttt{cb}data set. If
these variables are not present, no exposure results. The variables may
be present, but an individual person may still receive zero exposure
because not all samples of that food are contaminated, as indicated by
the \texttt{zeros.FV} value. The \texttt{nonzeros.FV} value is used to
determine the likelihood of contamination, with the residue variable
determining the amount found (when it is nonzero). The dietary exposure
is the product of the \texttt{consumption} as reported in the
\texttt{ftypes} input (in grams) and the \texttt{residue.FV} (in
micrograms of chemical per gram of food), summed over all food groups.
Three other vectors are returned (corresponding to dermal exposure,
inhalation exposure, and inhalation dose), but these are currently set
to zero for the dietary pathway. In principle, eating food could result
in dermal exposure (for finger foods), or inhalation (as foods may have
noticeable odors), but such exposures are not large enough to be of
concern at present.

\hypertarget{author-17}{%
\subsection{Author}\label{author-17}}

Kristin Isaacs, Graham Glen

\hypertarget{seealso-13}{%
\subsection{Seealso}\label{seealso-13}}

\texttt{{[}Diet\_diaries{]}(Diet\_diaries)}, \texttt{{[}run{]}(run)},
\texttt{{[}p.round{]}(p.round)},
\texttt{{[}generate.person.vars{]}(generate.person.vars)},
\texttt{{[}food.migration{]}(food.migration)}

\hypertarget{value-18}{%
\subsection{Value}\label{value-18}}

dietary The calculated quantity of chemical exposure from food residue
in grams per person per day.

\hypertarget{gen.factor.tables-gen.factor.tables}{%
\section{\texorpdfstring{\texttt{gen.factor.tables}:
gen.factor.tables}{gen.factor.tables: gen.factor.tables}}\label{gen.factor.tables-gen.factor.tables}}

\hypertarget{description-24}{%
\subsection{Description}\label{description-24}}

Constructs tables of non-media specific exposure factors for each
relevant combination of age, gender, and season. The input is from the
\texttt{{[}Exp\_actors{]}(Exp\_actors)} file (specified on the
\texttt{{[}Run{]}(Run)} file) after being read through the
\texttt{{[}read.exp.factors{]}(read.exp.factors)} function.

\hypertarget{usage-24}{%
\subsection{Usage}\label{usage-24}}

\begin{Shaded}
\begin{Highlighting}[]
\FunctionTok{gen.factor.tables}\NormalTok{(}\AttributeTok{ef =}\NormalTok{ exp.factors)}
\end{Highlighting}
\end{Shaded}

\hypertarget{arguments-21}{%
\subsection{Arguments}\label{arguments-21}}

\begin{itemize}
\tightlist
\item
  \texttt{ef}: A data set created internally using the
  \texttt{{[}run{]}(run)} function and the
  \texttt{{[}read.exp.factors{]}(read.exp.factors)} function to import
  the user specified \texttt{{[}Exp\_factors{]}(Exp\_factors)} input
  file. The data set contains the distributional parameters for the
  exposure factors. All of these variables may have age or
  gender-dependent distributions, although in the absence of data, many
  are assigned a single distribution from which all persons are sampled.
\end{itemize}

\hypertarget{author-18}{%
\subsection{Author}\label{author-18}}

Kristin Isaacs, Graham Glen

\hypertarget{seealso-14}{%
\subsection{Seealso}\label{seealso-14}}

\texttt{{[}Exp\_factors{]}(Exp\_factors)}, \texttt{{[}Run{]}(Run)},
\texttt{{[}run{]}(run)},
\texttt{{[}read.exp.factors{]}(read.exp.factors)}

\hypertarget{value-19}{%
\subsection{Value}\label{value-19}}

exp.gen Output consists of non-media specific exposure factors as a data
table. Depending on the user input, these generated exposure factors may
be gender and/or season specific. Each gender-season combination is
assigned a row number pointing to the appropriate distribution on the
output dataset from the
\texttt{{[}read.exp.factors{]}(read.exp.factors)} function. Thus,
\texttt{exp.gen} consists of 8 rows per variable (2 genders x 4
seasons), regardless of the number of different distributions used for
that variable. In addition, if the function runs successfully, the
following message is printed: ``General Factor Tables completed''

\hypertarget{get.fug.concs-get.fug.concs}{%
\section{\texorpdfstring{\texttt{get.fug.concs}:
get.fug.concs}{get.fug.concs: get.fug.concs}}\label{get.fug.concs-get.fug.concs}}

\hypertarget{description-25}{%
\subsection{Description}\label{description-25}}

Performs fugacity calculations to evaluate time-dependent chemical
flows.

\hypertarget{usage-25}{%
\subsection{Usage}\label{usage-25}}

\begin{Shaded}
\begin{Highlighting}[]
\FunctionTok{get.fug.concs}\NormalTok{(sdata, chem.data, x, cfug)}
\end{Highlighting}
\end{Shaded}

\hypertarget{arguments-22}{%
\subsection{Arguments}\label{arguments-22}}

\begin{itemize}
\tightlist
\item
  \texttt{sdata}: The chemical-scenario data specific to relevant
  combinations of chemical and scenario. Generated internally.
\item
  \texttt{chem.data}: The list of scenario-specific information for the
  chemicals being evaluated. Generated internally.
\item
  \texttt{x}: The output of the \texttt{{[}add.fugs{]}(add.fugs)}
  function.
\item
  \texttt{cfug}: The output of the \texttt{{[}chem\_fug{]}(chem\_fug)}
  function, a data set with the chemical specific parameters for each
  combination of chemical and simulated person. For each chemical, the
  chemical specific parameters assigned to a given person are randomly
  sampled from distributions on those parameters. These distributions
  are created from point estimates to reflect real-world uncertainty and
  variability.
\end{itemize}

\hypertarget{keyword-2}{%
\subsection{Keyword}\label{keyword-2}}

SHEDS

\hypertarget{keyword-3}{%
\subsection{Keyword}\label{keyword-3}}

kwd2

\hypertarget{details-17}{%
\subsection{Details}\label{details-17}}

This is one of two functions that perform fugacity calculations. This
one evaluates dynamic or time-dependent chemical First, a set of local
variables are determined for use in later calculations. These are a mix
of fixed and chemical-dependent variables (evaluated separately for each
person). Some variables, like chemical mass and app.rates, vary with
each source, so these calculations are repeated for each
source-scenario. Second, the eigenvalues and eigenvectors of the
jacobian matrix are calculated. Since the fugacity model has been
reduced to just two compartments (air and surface), the solutions can be
expressed analytically, and there is no explicit invocation of any
linear algebra routines that would normally be required. Third, the
variables composing the \texttt{concs} output are evaluated. The
variables \texttt{m.c.air} and \texttt{m.c.sur}are the time-constant
masses, while \texttt{m.t0.air} and \texttt{m.t0.sur} are the
time-dependent masses at t=0. The time-constant parts are zero here
because the permanent sources (i.e.~\texttt{c.src.air} and
\texttt{c.src.sur}) are assumed to be zero in these calculations. The
time-dependent masses are multiplied exponentially as a function of
time, and thus approach zero when enough time has passed.

\hypertarget{author-19}{%
\subsection{Author}\label{author-19}}

Kristin Isaacs, Graham Glen

\hypertarget{seealso-15}{%
\subsection{Seealso}\label{seealso-15}}

\texttt{{[}chem\_fugs{]}(chem\_fugs)},
\texttt{{[}Fugacity{]}(Fugacity)}, \texttt{{[}Run{]}(Run)},
\texttt{{[}run{]}(run)}, \texttt{{[}get.y0.concs{]}(get.y0.concs)}

\hypertarget{value-20}{%
\subsection{Value}\label{value-20}}

concs A data set containing calculated dynamic chemical flows for each
unique combination of simulated person and chemical.

\hypertarget{get.y0.concs-get.y0.concs}{%
\section{\texorpdfstring{\texttt{get.y0.concs}:
get.y0.concs}{get.y0.concs: get.y0.concs}}\label{get.y0.concs-get.y0.concs}}

\hypertarget{description-26}{%
\subsection{Description}\label{description-26}}

Performs fugacity calculations to evaluate constant (time-independent)
chemical flows, such as emissions from household articles.

\hypertarget{usage-26}{%
\subsection{Usage}\label{usage-26}}

\begin{Shaded}
\begin{Highlighting}[]
\FunctionTok{get.y0.concs}\NormalTok{(sdata, chem.data, pdmff, cfug)}
\end{Highlighting}
\end{Shaded}

\hypertarget{arguments-23}{%
\subsection{Arguments}\label{arguments-23}}

\begin{itemize}
\tightlist
\item
  \texttt{sdata}: The chemical-scenario data specific to relevant
  combinations of chemical and scenario. Generated internally.
\item
  \texttt{chem.data}: The list of scenario-specific information for the
  chemicals being evaluated. Generated internally.
\item
  \texttt{pdmff}: Output from the \texttt{{[}add.fugs{]}(add.fugs)}
  function. A data set containing values sampled from the distributions
  of each relevant variable in the \texttt{{[}Fugacity{]}(Fugacity)}
  input file for each theoretical person.
\item
  \texttt{cfug}: The output of the \texttt{{[}chem\_fug{]}(chem\_fug)}
  function, a data set with the chemical specific parameters for each
  combination of chemical and simulated person. For each chemical, the
  chemical specific parameters assigned to a given person are randomly
  sampled from distributions on those parameters. These distributions
  are created from point estimates to reflect real-world uncertainty and
  variability.
\end{itemize}

\hypertarget{keyword-4}{%
\subsection{Keyword}\label{keyword-4}}

SHEDS

\hypertarget{keyword-5}{%
\subsection{Keyword}\label{keyword-5}}

kwd2

\hypertarget{details-18}{%
\subsection{Details}\label{details-18}}

This function evaluates the chemical concentrations resulting from
constant source emissions. Thus, the function employs the steady state
solution to the fugacity equations. The basis for these calculations is
that the chemical sources are from articles, and are thus permanent and
unchanging. The chemical concentrations will therefore quickly adjust so
that the flows are balanced, and the concentrations remain fixed
thereafter.

\hypertarget{author-20}{%
\subsection{Author}\label{author-20}}

Kristin Isaacs, Graham Glen

\hypertarget{seealso-16}{%
\subsection{Seealso}\label{seealso-16}}

\texttt{{[}chem\_fugs{]}(chem\_fugs)},
\texttt{{[}Fugacity{]}(Fugacity)}, \texttt{{[}Run{]}(Run)},
\texttt{{[}run{]}(run)}, \texttt{{[}get.y0.concs{]}(get.y0.concs)}

\hypertarget{value-21}{%
\subsection{Value}\label{value-21}}

concs A data set containing calculated steady state chemical flows for
each unique combination of simulated person and chemical.

\hypertarget{indir.exposure-indir.exposure}{%
\section{\texorpdfstring{\texttt{indir.exposure}:
indir.exposure}{indir.exposure: indir.exposure}}\label{indir.exposure-indir.exposure}}

\hypertarget{description-27}{%
\subsection{Description}\label{description-27}}

Models the indirect exposure to chemicals in the home for each
theoretical person.

\hypertarget{usage-27}{%
\subsection{Usage}\label{usage-27}}

\begin{Shaded}
\begin{Highlighting}[]
\FunctionTok{indir.exposure}\NormalTok{(sd, cb, concs, chem.data)}
\end{Highlighting}
\end{Shaded}

\hypertarget{arguments-24}{%
\subsection{Arguments}\label{arguments-24}}

\begin{itemize}
\tightlist
\item
  \texttt{sd}: The chemical-scenario data specific to relevant
  combinations of chemical and scenario. Generated internally.
\item
  \texttt{cb}: A copy of the \texttt{base} data set output from the
  \texttt{{[}make.cbase{]}(make.cbase)} function, with columns added for
  exposure variables.
\item
  \texttt{concs}: The concentration of the chemical (in air and/or on
  surfaces) being released into the environment. Outpu of the
  \texttt{{[}get.fug.concs{]}(get.fug.concs)} function.
\item
  \texttt{chem.data}: The list of scenario-specific information for the
  chemicals being evaluated. Generated internally.
\end{itemize}

\hypertarget{details-19}{%
\subsection{Details}\label{details-19}}

Indirect exposure happens after a product is no longer being used or
applied, due to chemical lingering on various surfaces or in the air.
People who come along later may receive dermal or inhalation exposure
from residual chemical in the environment. SHEDS.HT currently has two
indirect exposure scenarios. One which applies to a one-time chemical
treatment applied to a house, and another which applies to continual
releases from articles. Both scenarios consist of two parts: the first
determines the appropriate air and surface concentrations. That code is
in the \texttt{{[}Fugacity{]}(Fugacity)} module. The second part is the
exposure calculation by the current function. Both types of indirect
exposure scenarios call this function. The surface and air
concentrations from the \texttt{{[}Fugacity{]}(Fugacity)} module are
premised on the product use actually occurring. Hence,
\texttt{{[}indir.exposure{]}(indir.exposure)} starts by multiplying
those concentrations by the \texttt{prevalence} (which is either 0 or 1,
evaluated separately for each person). For air, the exposure is the
average daily concentration, which is the event concentration multiplied
by the fraction of the day spent in that event. The inhaled dose is the
product of the exposure, the basal ventilation rate, and the PAI factor
(multiplier for the basal rate). A factor of 1E+06 converts the result
from grams per day to micrograms per day. Dermal exposure results from
skin contact with surfaces. The surface concentration (ug/cm2) is
multiplied by the fraction available for transfer (\texttt{avail.f},
unitless), the transfer coefficient (\texttt{dermal.tc}, in cm2/hr), and
the contact \texttt{duration} (hr/day). The result is the amount of
chemical transferred onto the skin (ug/day).

\hypertarget{author-21}{%
\subsection{Author}\label{author-21}}

Kristin Isaacs, Graham Glen

\hypertarget{seealso-17}{%
\subsection{Seealso}\label{seealso-17}}

\texttt{{[}Fugacity{]}(Fugacity)},
\texttt{{[}get.fug.concs{]}(get.fug.concs)},
\texttt{{[}make.cbase{]}(make.cbase)}, \texttt{{[}run{]}(run)}

\hypertarget{value-22}{%
\subsection{Value}\label{value-22}}

indir Indirect exposure to chemicals in the home in ug per person per
day.

\hypertarget{make.cbase-make.cbase}{%
\section{\texorpdfstring{\texttt{make.cbase}:
make.cbase}{make.cbase: make.cbase}}\label{make.cbase-make.cbase}}

\hypertarget{description-28}{%
\subsection{Description}\label{description-28}}

Extends the \texttt{base} data set, output from the
\texttt{{[}generate.person.vars{]}(generate.person.vars)} function, to
include a set of chemical-specific exposure variables. Currently, all
new variables are initialized to zero.

\hypertarget{usage-28}{%
\subsection{Usage}\label{usage-28}}

\begin{Shaded}
\begin{Highlighting}[]
\FunctionTok{make.cbase}\NormalTok{(base, chem)}
\end{Highlighting}
\end{Shaded}

\hypertarget{arguments-25}{%
\subsection{Arguments}\label{arguments-25}}

\begin{itemize}
\tightlist
\item
  \texttt{base}: Data set of all chemical-independent information needed
  for the exposure assessment. Each row corresponds to a simulated
  person. The variables consist of age, gender, weight, diet and
  activity diaries, food consumption, minutes in each micro, and the
  evaluation of the exposure factors for that person.
\item
  \texttt{chem}: List of the chemical(s) of interest, determined via the
  \texttt{chemical} input in the \texttt{{[}Run{]}(Run)} file. The
  exposure variables for the specified chemicals will be appended to the
  \texttt{base} data set.
\end{itemize}

\hypertarget{author-22}{%
\subsection{Author}\label{author-22}}

Kristin Isaacs, Graham Glen

\hypertarget{seealso-18}{%
\subsection{Seealso}\label{seealso-18}}

\texttt{{[}Run{]}(Run)}, \texttt{{[}run{]}(run)}

\hypertarget{value-23}{%
\subsection{Value}\label{value-23}}

The output is a data set cb consisting of the base data set, appended by
the following columns:

\begin{itemize}
\tightlist
\item
  chemrep: The unique ID corresponding to each chemical specified in the
  \texttt{chem} argument.
\item
  inhal.abs.f: For each person, the fraction of absorption of a given
  chemical when inhaled.
\item
  urine.f: For each person, the fraction of intake of a given chemical
  that is excreted through urine.
\item
  exp.inhal.tot: For each person, the total exposure of a given chemical
  via the direct inhalation scenario.
\item
  dose.inhal.tot: For each person, the corresponding dose for total
  exposure of a given chemical via the direct inhalation scenario.
\item
  exp.dermal.tot: For each person, the total exposure of a given
  chemical via the direct dermal application scenario.
\item
  exp.ingest.tot: For each person, the total exposure of a given
  chemical via the direct ingestion scenario.
\item
  exp.dietary.tot: For each person, the total exposure of a given
  chemical via the dietary scenario.
\item
  exp.migrat.tot: For each person, the total exposure of a given
  chemical via the migration scenario.
\item
  exp.nondiet.tot: For each person, the total exposure of a given
  chemical via the non-dietary food exposure scenario.
\item
  exp.ddd.tot: For each person, the total exposure of a given chemical
  via the down the drain scenario.
\item
  exp.window.tot: For each person, the total exposure of a given
  chemical via the out the window scenario.
\end{itemize}

\hypertarget{med.factor.tables-med.factor.tables}{%
\section{\texorpdfstring{\texttt{med.factor.tables}:
med.factor.tables}{med.factor.tables: med.factor.tables}}\label{med.factor.tables-med.factor.tables}}

\hypertarget{description-29}{%
\subsection{Description}\label{description-29}}

Constructs tables of media specific exposure factors for each relevant
combination of age, gender, and season, and each of three media specific
variables in the \texttt{ef} argument. These are avail.f, dermal.tc, and
om.ratio.

\hypertarget{usage-29}{%
\subsection{Usage}\label{usage-29}}

\begin{Shaded}
\begin{Highlighting}[]
\FunctionTok{med.factor.tables}\NormalTok{(ef, media.sur)}
\end{Highlighting}
\end{Shaded}

\hypertarget{arguments-26}{%
\subsection{Arguments}\label{arguments-26}}

\begin{itemize}
\tightlist
\item
  \texttt{ef}: A data set created internally using the
  \texttt{{[}run{]}(run)} function and the
  \texttt{{[}read.exp.factors{]}(read.exp.factors)} function to import
  the user specified \texttt{{[}Exp\_factors{]}(Exp\_factors)} input
  file. The data set contains the distributional parameters for the
  exposure factors. All of these variables may have age or
  gender-dependent distributions, although in the absence of data, many
  are assigned a single distribution from which all persons are sampled.
\item
  \texttt{media.sur}: A list of surface media. This data set is created
  internally by sub-setting the \texttt{{[}Media{]}(Media)} input file
  (read in with the \texttt{{[}read.media.file{]}(read.media.file)}
  function within the \texttt{{[}run{]}(run)} function) to extract only
  surface media.
\end{itemize}

\hypertarget{note-2}{%
\subsection{Note}\label{note-2}}

The first input argument to \texttt{med.factor.tables} (ef) is created
by reading in the \texttt{{[}Exp\_factors{]}(Exp\_factors)}input file
(specified on the \texttt{{[}Run{]}(Run)} file) with the
\texttt{{[}read.exp.factors{]}(read.exp.factors)} function within the
\texttt{{[}run{]}(run)} function. The \texttt{media.sur} input is
created by sub-setting the \texttt{{[}Media{]}(Media)} input file (read
in with the \texttt{{[}read.media.file{]}(read.media.file)} function
within the \texttt{{[}run{]}(run)} function).

\hypertarget{author-23}{%
\subsection{Author}\label{author-23}}

Kristin Isaacs, Graham Glen

\hypertarget{seealso-19}{%
\subsection{Seealso}\label{seealso-19}}

\texttt{{[}Media{]}(Media)}, \texttt{{[}Exp\_factors{]}(Exp\_factors)},
\texttt{{[}Run{]}(Run)}, \texttt{{[}run{]}(run)},
\texttt{{[}read.exp.factors{]}(read.exp.factors)},
\texttt{{[}read.media.file{]}(read.media.file)}

\hypertarget{value-24}{%
\subsection{Value}\label{value-24}}

exp.med Media specific exposure factors presented as a data table. The
present version of the model consists of only 3 such exposure factors,
but the code will accept more. Depending on the user input, these
generated exposure factors may be gender and/or season specific in
addition to media specific. The output consists of 24 rows per variable
(2 genders x 4 seasons x 3 media), when all ages share the same
distribution. If a variable has N age categories (each with its own
distribution) then there are (24 x N) rows for that variable. In
addition, if the function runs successfully, the following message is
printed: ``Media-specific Factor Tables completed''

\hypertarget{p.round-p.round}{%
\section{\texorpdfstring{\texttt{p.round}:
p.round}{p.round: p.round}}\label{p.round-p.round}}

\hypertarget{description-30}{%
\subsection{Description}\label{description-30}}

p.round performs stochastic or probabilistic rounding of non-integer
values.

\hypertarget{usage-30}{%
\subsection{Usage}\label{usage-30}}

\begin{Shaded}
\begin{Highlighting}[]
\FunctionTok{p.round}\NormalTok{(}\AttributeTok{x =} \ConstantTok{NULL}\NormalTok{, }\AttributeTok{q =} \ConstantTok{NA}\NormalTok{)}
\end{Highlighting}
\end{Shaded}

\hypertarget{arguments-27}{%
\subsection{Arguments}\label{arguments-27}}

\begin{itemize}
\tightlist
\item
  \texttt{x}: Required Default = none
\item
  \texttt{q}: Optional Default = none
\end{itemize}

\hypertarget{details-20}{%
\subsection{Details}\label{details-20}}

The input ``x'' is a vector of values to be rounded. Each value of x is
rounded independently, either up or down. The fractional value of x is
used to determine weights for rounding in each direction. For example,
if x=4.3, then it is rounded up to 5 with probability 0.3, else rounded
down to 4 (with probability 0.7). The p.round function always returns
integer values, and does not introduce bias. In the above example, if
the argument were 4.3 a large number of times, then the returned values
(all either 4 or 5) would average 4.3 with a small residual error which
approaches zero as ``n'' gets large.

\hypertarget{author-24}{%
\subsection{Author}\label{author-24}}

Kristin Isaacs, Graham Glen

\hypertarget{value-25}{%
\subsection{Value}\label{value-25}}

A vector of the same length as the input ``x'', containing all integer
values.

\hypertarget{post.exposure-post.exposure}{%
\section{\texorpdfstring{\texttt{post.exposure}:
post.exposure}{post.exposure: post.exposure}}\label{post.exposure-post.exposure}}

\hypertarget{description-31}{%
\subsection{Description}\label{description-31}}

Resolves the fate of chemical after initial exposure has occurred. This
involves parsing out the amount of chemical removed and amount
ultimately contributing to the exposure dose for each person.

\hypertarget{usage-31}{%
\subsection{Usage}\label{usage-31}}

\begin{Shaded}
\begin{Highlighting}[]
\FunctionTok{post.exposure}\NormalTok{(cb, cprops)}
\end{Highlighting}
\end{Shaded}

\hypertarget{arguments-28}{%
\subsection{Arguments}\label{arguments-28}}

\begin{itemize}
\tightlist
\item
  \texttt{cb}: A copy of the \texttt{base} data set output from the
  \texttt{{[}make.cbase{]}(make.cbase)} function, with columns added for
  exposure variables.
\item
  \texttt{cprops}: The chemical properties required for SHEDS-HT. The
  default file (the
  \texttt{{[}Chem\_props\ file{]}(Chem\_props\%20file)} read in by the
  \texttt{{[}read.chem.props{]}(read.chem.props)} function and modified
  before input into the current function) was prepared from publicly
  available databases using a custom program (not part of SHEDS-HT). The
  default file contains 7 numerical inputs per chemical, and the
  required properties are molecular weight (\texttt{MW}), vapor pressure
  (\texttt{VP.Pa}), solubility (\texttt{water.sol.mg.l}), octanol-water
  partition coefficient (\texttt{log.Kow}), air decay rate
  (\texttt{half.air.hr}), decay rate on surfaces
  (\texttt{half.sediment.hr}), and permeability coefficient
  (\texttt{Kp}).
\end{itemize}

\hypertarget{details-21}{%
\subsection{Details}\label{details-21}}

This function resolves the fate of chemical after initial exposure has
occurred. The same function applies to all exposure scenarios. The
dermal exposure is the most complicated, as there are five removal
methods. All five are randomly sampled. While the sum of the five means
is close to one, the sum of five random samples might not be, so these
samples are treated as fractions of their sum. The \texttt{rem.bath}
variable is either 0 or 1, so the bath removal term is either zero or
the sampled bath removal efficiency. The \texttt{rem.brush} term is also
simple. The handwashing (\texttt{rem.wash}) and hand-to-mouth
(\texttt{rem.hmouth}) transfer terms are non-linear, because higher
frequencies have less chemical available for removal on each repetition.
The algorithms in place were fitted to output from SHEDS-Multimedia,
which were summed to daily totals. The final removal term is dermal
absorption (\texttt{rem.absorb}). The base value is multiplied by the
\texttt{Kp} factor from the \texttt{cprops} argument, and divided by the
value for permethrin, as the values from SHEDS-Multimedia were based on
a permethrin run. The five terms are evaluated separately for each
person, as is their sum. Each is then converted to a fraction of the
whole. The fractions may be quite different from one person to another.
For one, perhaps 80\% of the dermal loading is removed by a bath/shower,
while for another person it is 0\% because they did not take one. For
the latter person, other four removal terms are (on average) five times
larger than for the former person, because together they account for
100\% of the removal, instead of just 20\%. The rest of the
\texttt{post.exposure} function is mostly a matter of bookkeeping. The
hand-to-mouth dermal removal term becomes an ingestion exposure term.
Summing exposures across routes is dubious, in part because inhalation
exposures use different units from the others, and because much of the
dermal exposure never enters the body. In addition, summing dermal and
ingestion exposures may double-count the hand-to-mouth term. However,
intake dose may be summed. In SHEDS, ``intake dose'' is the sum of the
inhaled dose (which is the amount of chemical entering the lungs in
ug/day), the ingestion exposure (which is the amount entering the GI
tract in ug/day), and the dermal absorption (the amount penetrating into
or through the skin, so it cannot otherwise be removed, in ug/day). The
``absorbed dose'' is also calculated: for dermal it is the same as the
intake dose, but for ingestion and inhalation there is another
absorption factor, which was set on the exposure factors input file. An
estimate of the chemical in urine (in ug/day) is made. Both the intake
dose and the absorbed dose are reported in both (ug/day) and in
(mg/kg/day). Note that the latter requires the body weights of each
individual. These cannot be obtained from the former just by knowing the
average body weight in SHEDS. The above variables for each simulated
person are written to the \texttt{fexp} object (the name stands for
``final exposure''). Each chemical writes over the previous
\texttt{fexp}, so the data must first be summarized and written to an
output file.

\hypertarget{author-25}{%
\subsection{Author}\label{author-25}}

Kristin Isaacs, Graham Glen

\hypertarget{seealso-20}{%
\subsection{Seealso}\label{seealso-20}}

\texttt{{[}make.cbase{]}(make.cbase)},
\texttt{{[}Chemprops\_small{]}(Chemprops\_small)}

\hypertarget{value-26}{%
\subsection{Value}\label{value-26}}

fexp Absorbed dose and intake dose of a given chemical for each
theoretical person being modeled for all exposure scenarios.

\hypertarget{puc.boxplot-puc.boxplot-is-a-ggplot-boxplot-of-one-or-multiple-srcmeans-variables.-the-y-axis-for-each-row-is-the-selected-output-variable-and-the-x-axis-is-higher-level-puc.}{%
\section{\texorpdfstring{\texttt{puc.boxplot}: puc.boxplot is a ggplot
boxplot of one or multiple srcMeans variables. The y-axis for each row
is the selected output variable and the x-axis is higher-level
PUC.}{puc.boxplot: puc.boxplot is a ggplot boxplot of one or multiple srcMeans variables. The y-axis for each row is the selected output variable and the x-axis is higher-level PUC.}}\label{puc.boxplot-puc.boxplot-is-a-ggplot-boxplot-of-one-or-multiple-srcmeans-variables.-the-y-axis-for-each-row-is-the-selected-output-variable-and-the-x-axis-is-higher-level-puc.}}

User can select multiple different output variables using the argument
output.variables. Data comes from srcMeans output file, so data points
are mean output.variable. The user can plot multiple different pathways
and each row of boxplots is a different variable

\hypertarget{description-32}{%
\subsection{Description}\label{description-32}}

puc.boxplot is a ggplot boxplot of one or multiple srcMeans variables.
The y-axis for each row is the selected output variable and the x-axis
is higher-level PUC. User can select multiple different output variables
using the argument output.variables. Data comes from srcMeans output
file, so data points are mean output.variable. The user can plot
multiple different pathways and each row of boxplots is a different
variable

\hypertarget{usage-32}{%
\subsection{Usage}\label{usage-32}}

\begin{Shaded}
\begin{Highlighting}[]
\FunctionTok{puc.boxplot}\NormalTok{(}
\NormalTok{  run.name,}
  \AttributeTok{output.variables =} \FunctionTok{c}\NormalTok{(}\StringTok{"exp.dermal"}\NormalTok{, }\StringTok{"exp.ingest"}\NormalTok{, }\StringTok{"exp.inhal"}\NormalTok{)}
\NormalTok{)}
\end{Highlighting}
\end{Shaded}

\hypertarget{arguments-29}{%
\subsection{Arguments}\label{arguments-29}}

\begin{itemize}
\tightlist
\item
  \texttt{run.name}: default = name of last run (in current R session)
\item
  \texttt{output.variables}: argument is a concatenated vector of
  character strings. c(``exp.dermal'', ``exp.ingest'', ``exp.inhal'')
\end{itemize}

\hypertarget{details-22}{%
\subsection{Details}\label{details-22}}

A SHEDS run creates an output file for each chemical. This plot pulls
from the srcMeans.csv for each chemical in the output file. The y axis
is log transformed and a constant of 1e-10 added to y variable for log
transformation of 0's.

\hypertarget{value-27}{%
\subsection{Value}\label{value-27}}

a boxplot of output.variable vs higher level PUC

\hypertarget{puc.rank.plot-puc.rank.plot-is-a-ggplot-scatter-plot-with-desired-srcmeans-variable-on-the-y-axis-and-chemical-rank-on-the-x-axis.}{%
\section{\texorpdfstring{\texttt{puc.rank.plot}: puc.rank.plot is a
ggplot scatter plot with desired srcMeans variable on the y axis and
chemical rank on the x
axis.}{puc.rank.plot: puc.rank.plot is a ggplot scatter plot with desired srcMeans variable on the y axis and chemical rank on the x axis.}}\label{puc.rank.plot-puc.rank.plot-is-a-ggplot-scatter-plot-with-desired-srcmeans-variable-on-the-y-axis-and-chemical-rank-on-the-x-axis.}}

The dependent variable on the y axis can be changed using
output.variable argument, but only one output variable can be plotted at
a time. User can either plot the individual product or the higher level
puc using the combine argument. Data comes from all\_srcMeans output
file, so data points are the mean for that output varaible.

\hypertarget{description-33}{%
\subsection{Description}\label{description-33}}

puc.rank.plot is a ggplot scatter plot with desired srcMeans variable on
the y axis and chemical rank on the x axis. The dependent variable on
the y axis can be changed using output.variable argument, but only one
output variable can be plotted at a time. User can either plot the
individual product or the higher level puc using the combine argument.
Data comes from all\_srcMeans output file, so data points are the mean
for that output varaible.

\hypertarget{usage-33}{%
\subsection{Usage}\label{usage-33}}

\begin{Shaded}
\begin{Highlighting}[]
\FunctionTok{puc.rank.plot}\NormalTok{(run.name, }\AttributeTok{output.variable =} \StringTok{"exp.dermal"}\NormalTok{, }\AttributeTok{combine =} \ConstantTok{FALSE}\NormalTok{)}
\end{Highlighting}
\end{Shaded}

\hypertarget{arguments-30}{%
\subsection{Arguments}\label{arguments-30}}

\begin{itemize}
\tightlist
\item
  \texttt{run.name}: default = name of last run (in current R session)
\item
  \texttt{output.variable}: argument is a single character string.
  default = ``exp.dermal''.
\item
  \texttt{combine}: true or false statement that dictates whether to
  plot each product or combine them so that each data point is a higher
  level product. default = FALSE
\end{itemize}

\hypertarget{details-23}{%
\subsection{Details}\label{details-23}}

A SHEDS run creates an output file for each chemical. This plot pulls
from the srcMeans.csv for each chemical in the output file. Chemicals
are ranked in ascending order of output.variable on the x axis. The y
axis is log transformed and a constant of 1e-10 added to y variable for
log transformation of 0's. The plot colors points by higher level PUC.

\hypertarget{value-28}{%
\subsection{Value}\label{value-28}}

a plot of output.variable vs chemical rank with a legend for color and
shape

\hypertarget{quantiles-quantiles}{%
\section{\texorpdfstring{\texttt{quantiles}:
quantiles}{quantiles: quantiles}}\label{quantiles-quantiles}}

\hypertarget{description-34}{%
\subsection{Description}\label{description-34}}

This function is similar to the built-in R function ``quantile'', but it
returns a list of pre-selected quantiles of the set of values in the
vector ``x''.

\hypertarget{usage-34}{%
\subsection{Usage}\label{usage-34}}

\begin{Shaded}
\begin{Highlighting}[]
\FunctionTok{quantiles}\NormalTok{(x)}
\end{Highlighting}
\end{Shaded}

\hypertarget{arguments-31}{%
\subsection{Arguments}\label{arguments-31}}

\begin{itemize}
\tightlist
\item
  \texttt{x}: Default = none
\end{itemize}

\hypertarget{details-24}{%
\subsection{Details}\label{details-24}}

This function is similar to the built-in R function ``quantile'', but it
returns a list of pre-selected quantiles of the set of values in the
vector ``x''. Specifically, it returns all the following quantiles:
.005, .01, .025, .05, .1, .15, .2, .25, .3, .4, .5, .6, .7, .75, .8,
.85, .9, .95, .975, .99, .995

\hypertarget{author-26}{%
\subsection{Author}\label{author-26}}

Kristin Isaacs, Graham Glen

\hypertarget{value-29}{%
\subsection{Value}\label{value-29}}

This function is used to construct the tables in the ``Allstats'' output
files.

\hypertarget{read.act.diaries-read.act.diaries}{%
\section{\texorpdfstring{\texttt{read.act.diaries}:
read.act.diaries}{read.act.diaries: read.act.diaries}}\label{read.act.diaries-read.act.diaries}}

\hypertarget{description-35}{%
\subsection{Description}\label{description-35}}

Read.act.diaries reads human activity diaries from the .csv file
indicated by ``filename''. ``Specs'' contains the run specifications
from the run.file, and is used to subset the activity diaries by age,
gender, or season, if requested.

\hypertarget{usage-35}{%
\subsection{Usage}\label{usage-35}}

\begin{Shaded}
\begin{Highlighting}[]
\FunctionTok{read.act.diaries}\NormalTok{(filename, specs)}
\end{Highlighting}
\end{Shaded}

\hypertarget{read.chem.props-read.chem.props}{%
\section{\texorpdfstring{\texttt{read.chem.props}:
read.chem.props}{read.chem.props: read.chem.props}}\label{read.chem.props-read.chem.props}}

\hypertarget{description-36}{%
\subsection{Description}\label{description-36}}

Read.chem.props reads the chemical properties from the .csv file
indicated by ``filename''. ``Specs'' contains the run specifications
from the run.file, and is used to subset the chemicals by the list
provided on the run.file. If no such list was given, then all chemicals
are kept.

\hypertarget{usage-36}{%
\subsection{Usage}\label{usage-36}}

\begin{Shaded}
\begin{Highlighting}[]
\FunctionTok{read.chem.props}\NormalTok{(filename, specs)}
\end{Highlighting}
\end{Shaded}

\hypertarget{read.diet.diaries-read.diet.diaries}{%
\section{\texorpdfstring{\texttt{read.diet.diaries}:
read.diet.diaries}{read.diet.diaries: read.diet.diaries}}\label{read.diet.diaries-read.diet.diaries}}

\hypertarget{description-37}{%
\subsection{Description}\label{description-37}}

Read.diet.diaries reads the food diaries from the .csv file indicated by
``filename''. ``Specs'' contains the run specifications from the
run.file, and is used to subset the activity diaries by age and/or
gender.

\hypertarget{usage-37}{%
\subsection{Usage}\label{usage-37}}

\begin{Shaded}
\begin{Highlighting}[]
\FunctionTok{read.diet.diaries}\NormalTok{(filename, specs)}
\end{Highlighting}
\end{Shaded}

\hypertarget{read.exp.factors-read.exp.factors}{%
\section{\texorpdfstring{\texttt{read.exp.factors}:
read.exp.factors}{read.exp.factors: read.exp.factors}}\label{read.exp.factors-read.exp.factors}}

\hypertarget{description-38}{%
\subsection{Description}\label{description-38}}

Read.exp.factors reads the exposure factors from the .csv file indicated
by ``filename''.

\hypertarget{usage-38}{%
\subsection{Usage}\label{usage-38}}

\begin{Shaded}
\begin{Highlighting}[]
\FunctionTok{read.exp.factors}\NormalTok{(filename)}
\end{Highlighting}
\end{Shaded}

\hypertarget{read.fug.inputs-read.fug.inputs}{%
\section{\texorpdfstring{\texttt{read.fug.inputs}:
read.fug.inputs}{read.fug.inputs: read.fug.inputs}}\label{read.fug.inputs-read.fug.inputs}}

\hypertarget{description-39}{%
\subsection{Description}\label{description-39}}

Read.fug.inputs reads the non-chemical dependent inputs for fugacity
modeling in SHEDS from the .csv file indicated by ``filename''.

\hypertarget{usage-39}{%
\subsection{Usage}\label{usage-39}}

\begin{Shaded}
\begin{Highlighting}[]
\FunctionTok{read.fug.inputs}\NormalTok{(filename)}
\end{Highlighting}
\end{Shaded}

\hypertarget{read.media.file-read.media.file}{%
\section{\texorpdfstring{\texttt{read.media.file}:
read.media.file}{read.media.file: read.media.file}}\label{read.media.file-read.media.file}}

\hypertarget{description-40}{%
\subsection{Description}\label{description-40}}

Read.media.file reads the names, properties, and associated
microenvironments for each of the potential exposure media in SHEDS.

\hypertarget{usage-40}{%
\subsection{Usage}\label{usage-40}}

\begin{Shaded}
\begin{Highlighting}[]
\FunctionTok{read.media.file}\NormalTok{(filename)}
\end{Highlighting}
\end{Shaded}

\hypertarget{read.phys.file-read.phys.file}{%
\section{\texorpdfstring{\texttt{read.phys.file}:
read.phys.file}{read.phys.file: read.phys.file}}\label{read.phys.file-read.phys.file}}

\hypertarget{description-41}{%
\subsection{Description}\label{description-41}}

Read.phys.file reads the physiology data from the .csv file indicated by
``filename''.

\hypertarget{usage-41}{%
\subsection{Usage}\label{usage-41}}

\begin{Shaded}
\begin{Highlighting}[]
\FunctionTok{read.phys.file}\NormalTok{(filename)}
\end{Highlighting}
\end{Shaded}

\hypertarget{read.pop.file-read.pop.file}{%
\section{\texorpdfstring{\texttt{read.pop.file}:
read.pop.file}{read.pop.file: read.pop.file}}\label{read.pop.file-read.pop.file}}

\hypertarget{description-42}{%
\subsection{Description}\label{description-42}}

Read.pop.file reads the population data from the .csv file indicated by
``filename''.

\hypertarget{usage-42}{%
\subsection{Usage}\label{usage-42}}

\begin{Shaded}
\begin{Highlighting}[]
\FunctionTok{read.pop.file}\NormalTok{(filename, specs)}
\end{Highlighting}
\end{Shaded}

\hypertarget{read.run.file-read.run.file}{%
\section{\texorpdfstring{\texttt{read.run.file}:
read.run.file}{read.run.file: read.run.file}}\label{read.run.file-read.run.file}}

\hypertarget{description-43}{%
\subsection{Description}\label{description-43}}

Each SHEDS run has its own ``run.file'' that the user prepares before
the run. This file contains all the settings and file references needed
for the run. Read.run.file occurs at the start of each SHEDS run. The
contents of the run.file are examined, and stored in the R object
``specs''.

\hypertarget{usage-43}{%
\subsection{Usage}\label{usage-43}}

\begin{Shaded}
\begin{Highlighting}[]
\FunctionTok{read.run.file}\NormalTok{(}\AttributeTok{run.file =} \StringTok{"run\_test.txt"}\NormalTok{)}
\end{Highlighting}
\end{Shaded}

\hypertarget{read.source.chem.file-read.source.chem.file}{%
\section{\texorpdfstring{\texttt{read.source.chem.file}:
read.source.chem.file}{read.source.chem.file: read.source.chem.file}}\label{read.source.chem.file-read.source.chem.file}}

\hypertarget{description-44}{%
\subsection{Description}\label{description-44}}

Read.source.chem.file reads the distributions that are specific to
combinations of source and chemical from the indicated .csv file.

\hypertarget{usage-44}{%
\subsection{Usage}\label{usage-44}}

\begin{Shaded}
\begin{Highlighting}[]
\FunctionTok{read.source.chem.file}\NormalTok{(filename, scenSrc, specs)}
\end{Highlighting}
\end{Shaded}

\hypertarget{read.source.scen.file-read.source.scen.file}{%
\section{\texorpdfstring{\texttt{read.source.scen.file}:
read.source.scen.file}{read.source.scen.file: read.source.scen.file}}\label{read.source.scen.file-read.source.scen.file}}

\hypertarget{description-45}{%
\subsection{Description}\label{description-45}}

Read.source.scen.file reads the list of active exposure scenarios for
each potential source of chemical from the .csv file indicated by
``filename''.

\hypertarget{usage-45}{%
\subsection{Usage}\label{usage-45}}

\begin{Shaded}
\begin{Highlighting}[]
\FunctionTok{read.source.scen.file}\NormalTok{(filename)}
\end{Highlighting}
\end{Shaded}

\hypertarget{read.source.vars.file-read.source.vars.file}{%
\section{\texorpdfstring{\texttt{read.source.vars.file}:
read.source.vars.file}{read.source.vars.file: read.source.vars.file}}\label{read.source.vars.file-read.source.vars.file}}

\hypertarget{description-46}{%
\subsection{Description}\label{description-46}}

Read.source.vars.file reads the distributions that are specific to
combinations of source and chemical from the indicated .csv file.

\hypertarget{usage-46}{%
\subsection{Usage}\label{usage-46}}

\begin{Shaded}
\begin{Highlighting}[]
\FunctionTok{read.source.vars.file}\NormalTok{(filename, src.scen)}
\end{Highlighting}
\end{Shaded}

\hypertarget{run-run}{%
\section{\texorpdfstring{\texttt{run}: run}{run: run}}\label{run-run}}

\hypertarget{description-47}{%
\subsection{Description}\label{description-47}}

Function to call the \texttt{{[}Run{]}(Run)} txt file, which consists of
user-defined parameters and calls to input files required to initialize
a SHEDS.HT run.

\hypertarget{usage-47}{%
\subsection{Usage}\label{usage-47}}

\begin{Shaded}
\begin{Highlighting}[]
\FunctionTok{run}\NormalTok{(}\AttributeTok{run.file =} \StringTok{""}\NormalTok{, }\AttributeTok{wd =} \StringTok{""}\NormalTok{)}
\end{Highlighting}
\end{Shaded}

\hypertarget{arguments-32}{%
\subsection{Arguments}\label{arguments-32}}

\begin{itemize}
\tightlist
\item
  \texttt{run.file}: The name of the run file to be used for a given
  run. Many different ``Run'''' files may be set up for special
  purposes.The one being invoked must be present in the inputs folder.
\item
  \texttt{wd}: The user's working directory. The working directory
  should contain an inputs folder, containing all necessary SHEDS.HT
  input files. The wd value should be set once for each SHEDS.HT
  installation by replacing the default value in the definition of
  run().
\end{itemize}

\hypertarget{details-25}{%
\subsection{Details}\label{details-25}}

This function is used to provide R with information necessary to produce
a SHEDS-HT run and to initialize necessary parameters. No values are
returned. In order to produce a successful run, all necessary input
files should be stored in the working directory specified by the wd
argument.

\hypertarget{author-27}{%
\subsection{Author}\label{author-27}}

Kristin Isaacs, Graham Glen

\hypertarget{seealso-21}{%
\subsection{Seealso}\label{seealso-21}}

\texttt{{[}Run.txt{]}(Run.txt)},
\texttt{{[}read.run.file{]}(read.run.file)}

\hypertarget{value-30}{%
\subsection{Value}\label{value-30}}

No variable will be returned.

\hypertarget{scen.factor.indices-scen.factor.indices}{%
\section{\texorpdfstring{\texttt{scen.factor.indices}:
scen.factor.indices}{scen.factor.indices: scen.factor.indices}}\label{scen.factor.indices-scen.factor.indices}}

\hypertarget{description-48}{%
\subsection{Description}\label{description-48}}

Constructs scenario-specific exposure factors for each relevant
combination of age and gender. The input is derived internally from the
\texttt{{[}Source\_vars{]}(Source\_vars)} file specified on the
\texttt{{[}Run{]}(Run)} file.

\hypertarget{usage-48}{%
\subsection{Usage}\label{usage-48}}

\begin{Shaded}
\begin{Highlighting}[]
\FunctionTok{scen.factor.indices}\NormalTok{(sdat, expgen)}
\end{Highlighting}
\end{Shaded}

\hypertarget{arguments-33}{%
\subsection{Arguments}\label{arguments-33}}

\begin{itemize}
\tightlist
\item
  \texttt{sdat}: The chemical-scenario data specific to a given
  combination of chemical and scenario.
\item
  \texttt{expgen}: Non-media specific exposure factors as a data table.
  Output from \texttt{{[}gen.factor.tables{]}(gen.factor.tables)}
\end{itemize}

\hypertarget{details-26}{%
\subsection{Details}\label{details-26}}

The constructed factors are not media-specific, although they are
scenario-specific, and most scenarios include just one surface medium.
This function evaluats the scenario-specific exposure factors separately
from the other factors because these scenario-specific factors may
change with every chemical and scenario, whereas the other factors
remain the same across chemicals and scenarios for each person.

\hypertarget{author-28}{%
\subsection{Author}\label{author-28}}

Kristin Isaacs, Graham Glen

\hypertarget{value-31}{%
\subsection{Value}\label{value-31}}

Returns indices of scenario-specific exposure factors for each relevant
age and gender combination. The output is only generated internally.

\hypertarget{select.people-select.people}{%
\section{\texorpdfstring{\texttt{select.people}:
select.people}{select.people: select.people}}\label{select.people-select.people}}

\hypertarget{description-49}{%
\subsection{Description}\label{description-49}}

Assigns demographic and physiological variables to each theoretical
person to be modeled.

\hypertarget{usage-49}{%
\subsection{Usage}\label{usage-49}}

\begin{Shaded}
\begin{Highlighting}[]
\CommentTok{\# S3 method for people}
\FunctionTok{select}\NormalTok{(n, pop, py, act.p, diet.p, act.d, diet.d, specs)}
\end{Highlighting}
\end{Shaded}

\hypertarget{arguments-34}{%
\subsection{Arguments}\label{arguments-34}}

\begin{itemize}
\tightlist
\item
  \texttt{n}: Number of persons.
\item
  \texttt{pop}: The population input; output of
  \texttt{{[}read.pop.file{]}(read.pop.file)} function. Contains counts
  by gender and each year of age from the 2000 U.S. census. When a large
  age range is modeled, this ensures that SHEDS chooses age and gender
  with the correct overall probability. \#'
\item
  \texttt{py}: Regression parameters on the three physiological
  variables of interest (weight, height, and body mass index) for
  various age and gender groups. Output of the
  \texttt{{[}read.phys.file{]}(read.phys.file)} function.
\item
  \texttt{act.p}: Activity diary pools; output of the
  \texttt{{[}act.diary.pools{]}(act.diary.pools)} function. Each element
  in this input is a list of acceptable activity diary numbers for each
  year of age, for each gender, weekend, and season combination.
\item
  \texttt{diet.p}: Dietary diary pools; output of the
  \texttt{{[}diet.diary.pools{]}(diet.diary.pools)} function. Each
  element in this input is a list of acceptable activity diary numbers
  for each year of age, for each gender, weekend, and season
  combination.
\item
  \texttt{act.d}: Activity diaries, which indicate the amount of time
  and level of metabolic activity in various `micros'. Each line of data
  represents one person-day (24 hours). Output of the
  \texttt{{[}read.act.diaries{]}(read.act.diaries)} function.
\item
  \texttt{diet.d}: Daily diaries of dietary consumption by food group;
  output of the \texttt{{[}read.diet.diaries{]}(read.diet.diaries)}
  function. Each line represents one person-day, with demographic
  variables followed by amounts (in grams/day) for a list of food types
  indicated by a short abbreviation on the header line.
\item
  \texttt{specs}: Output of the
  \texttt{{[}read.run.file{]}(read.run.file)} function, which can be
  modified by the \texttt{{[}update.specs{]}(update.specs)} function
  before input into \texttt{{[}select.people{]}(select.people)}.
\end{itemize}

\hypertarget{details-27}{%
\subsection{Details}\label{details-27}}

This is the first real step in the modeling process. It first fills an
array \texttt{q} with uniform random numbers, with ten columns because
there are 10 random variables defined by this function. There number of
rows correspond to the number of persons, capped at the
\texttt{set.size} specified in the \texttt{{[}Run{]}(Run)} input file
(typically 5000). Gender is selected from a discrete (binomial)
distribution where the counts of males and females in the study age
range determines the gender probabilities. Age is tabulated next,
separately for each gender. The counts by year of age are chosen for the
appropriate gender and used as selection weights. Season is assigned
randomly (equal weights) using those specified in the
\texttt{{[}Run{]}(Run)}input file. \texttt{Weekend} is set to one or
zero, with a chance of 2/7 for the former. The next block of code
assigns physiological variables. Weight is lognormal in SHEDS, so a
normal is sampled first and then \texttt{exp()} is applied. This means
that the weight parameters refer to the properties of log(weight), which
were fit by linear regression. The basal metabolic rate (bmr), is
calculated by regression. A minimum bmr is set to prevent extreme cases
from becoming zero or negative. The alveolar breathing ventilation rate
corresponding to bmr is also calculated. The SHEDS logic sets activities
in each micro to be a multiple of these rates, with outdoor rates higher
than indoor, and indoor rates higher than sleep rates. This calculated
activities affect the inhaled dose. The skin surface area is calculated
using regressions based on height and weight for 3 age ranges. The next
step is to assign diaries. Here, a FOR loop over n persons (n rows) is
used to assign appropriate diet and activity diary pools to each person.
An empirical distribution is created, consisting of the list of diary
numbers for each pool. The final step is to retrieve the actual data
from the chosen activity and diet diaries, and the result becomes
\texttt{pd}.

\hypertarget{author-29}{%
\subsection{Author}\label{author-29}}

Kristin Isaacs, Graham Glen

\hypertarget{seealso-22}{%
\subsection{Seealso}\label{seealso-22}}

\texttt{{[}run{]}(run)}, \texttt{{[}read.pop.file{]}(read.pop.file)},
\texttt{{[}read.phys.file{]}(read.phys.file)},
\texttt{{[}act.diary.pools{]}(act.diary.pools)},
\texttt{{[}diet.diary.pools{]}(diet.diary.pools)},
\texttt{{[}read.act.diaries{]}(read.act.diaries)},
\texttt{{[}read.diet.diaries{]}(read.diet.diaries)},
\texttt{{[}update.specs{]}(update.specs)}

\hypertarget{value-32}{%
\subsection{Value}\label{value-32}}

pd A dataframe of ``person-demographics'': assigned demographic and
physiological parameters for each theoretical person modeled in SHEDS.HT

\hypertarget{set.pars-set.pars}{%
\section{\texorpdfstring{\texttt{set.pars}:
set.pars}{set.pars: set.pars}}\label{set.pars-set.pars}}

Adjusts certain SHEDS input distribution types to conform with the
requirements of the \texttt{{[}distrib{]}(distrib)} function.

\hypertarget{description-50}{%
\subsection{Description}\label{description-50}}

set.pars Adjusts certain SHEDS input distribution types to conform with
the requirements of the \texttt{{[}distrib{]}(distrib)} function.

\hypertarget{usage-50}{%
\subsection{Usage}\label{usage-50}}

\begin{Shaded}
\begin{Highlighting}[]
\FunctionTok{set.pars}\NormalTok{(vars)}
\end{Highlighting}
\end{Shaded}

\hypertarget{arguments-35}{%
\subsection{Arguments}\label{arguments-35}}

\begin{itemize}
\tightlist
\item
  \texttt{vars}: Output of the
  \texttt{{[}read.source.chem.file{]}(read.source.chem.file)} or the
  \texttt{{[}read.source.vars.file{]}(read.source.vars.file)} functions.
\end{itemize}

\hypertarget{details-28}{%
\subsection{Details}\label{details-28}}

Set.pars adjusts certain SHEDS input distribution types to conform with
the requirements of the Distrib function. The lognormal parameters are
changed from arithmetic mean and standard deviation to geometric mean
and geometric standard deviation. For the normal distribution, the
standard deviation is computed from the mean and coefficient of
variation (CV). For user prevalence, the input may be specified either
as a point value (indicating probability) or as a Bernoulli
distribution. If the former is used, it is converted to the latter.

\hypertarget{author-30}{%
\subsection{Author}\label{author-30}}

Kristin Isaacs, Graham Glen

\hypertarget{seealso-23}{%
\subsection{Seealso}\label{seealso-23}}

\texttt{{[}run{]}(run)},
\texttt{{[}read.source.chem.file{]}(read.source.chem.file)},
\texttt{{[}read.source.vars.file{]}(read.source.vars.file)}

\hypertarget{value-33}{%
\subsection{Value}\label{value-33}}

v The modified version of the input vars.

\hypertarget{setup-setup}{%
\section{\texorpdfstring{\texttt{setup}:
setup}{setup: setup}}\label{setup-setup}}

\hypertarget{description-51}{%
\subsection{Description}\label{description-51}}

Loads required R packages and sources the modules needed to perform a
SHEDS.HT run. The user might need to download the packages if they are
not already present (see Dependencies).

\hypertarget{usage-51}{%
\subsection{Usage}\label{usage-51}}

\begin{Shaded}
\begin{Highlighting}[]
\FunctionTok{setup}\NormalTok{(}\AttributeTok{wd =} \StringTok{""}\NormalTok{)}
\end{Highlighting}
\end{Shaded}

\hypertarget{arguments-36}{%
\subsection{Arguments}\label{arguments-36}}

\begin{itemize}
\tightlist
\item
  \texttt{wd}: The User's working directory (set in the inputs to the
  \texttt{{[}run{]}(run)} function). Should contain all necessary
  SHEDS.HT inputs in an /Input folder.
\end{itemize}

\hypertarget{note-3}{%
\subsection{Note}\label{note-3}}

Requires: \texttt{{[}data.table{]}(data.table)},
\texttt{{[}plyr{]}(plyr)}, \texttt{{[}stringr{]}(stringr)},
\texttt{{[}ggplot2{]}(ggplot2)}

\hypertarget{author-31}{%
\subsection{Author}\label{author-31}}

Kristin Isaacs, Graham Glen

\hypertarget{seealso-24}{%
\subsection{Seealso}\label{seealso-24}}

\texttt{{[}run{]}(run)}

\hypertarget{value-34}{%
\subsection{Value}\label{value-34}}

No values returned.

\hypertarget{summarize.chemical-summarize.chemcial}{%
\section{\texorpdfstring{\texttt{summarize.chemical}:
summarize.chemcial}{summarize.chemical: summarize.chemcial}}\label{summarize.chemical-summarize.chemcial}}

\hypertarget{description-52}{%
\subsection{Description}\label{description-52}}

Summarize.chemical writes a .csv file containing a summary of the
exposure and dose results from the object ``x''.

\hypertarget{usage-52}{%
\subsection{Usage}\label{usage-52}}

\begin{Shaded}
\begin{Highlighting}[]
\FunctionTok{summarize.chemical}\NormalTok{(x, c, chem, chemical, set, sets, specs)}
\end{Highlighting}
\end{Shaded}

\hypertarget{arguments-37}{%
\subsection{Arguments}\label{arguments-37}}

\begin{itemize}
\tightlist
\item
  \texttt{x}: Default = none Exposure data set
\item
  \texttt{c}: Default = none Index \# for chemical
\item
  \texttt{chem}: Default = none CAS for chemical
\item
  \texttt{chemical}: Default = none Full chemical name
\item
  \texttt{set}: Default = none Index \# for set of simulated persons
\item
  \texttt{sets}: Default = none Total \# of sets in this SHEDS run
\item
  \texttt{specs}: Default = none List of settings from the ``run'' input
  file
\end{itemize}

\hypertarget{details-29}{%
\subsection{Details}\label{details-29}}

Tables are produced for each of the following cohorts: males, females,
females ages 16-49, age 0-5, age 6-11, age 12-19, age 20-65, age 66+,
and a table for all persons. The variables that are summarized are:
dermal exposure, ingestion exposure, inhalation exposure, inhaled dose,
intake dose, dermal absorption, ingestion absorption, inhalation
absorption, total absorption in micrograms per day, total absorption in
milligrams per kilogram per day, and chemical mass down the drain. In
each table, the following statistics are computed across the appropriate
subpopulation: mean, standard deviation, and quantiles .005, .01, .025,
.05, .1, .15, .2, .25, .3, .4, .5, .6, .7, .75, .8, .85, .9, .95, .975,
.99, .995.

\hypertarget{author-32}{%
\subsection{Author}\label{author-32}}

Kristin Isaacs, Graham Glen

\hypertarget{value-35}{%
\subsection{Value}\label{value-35}}

If ``x'' is a single set of data (that is, if the argument ``set'' is
between 1 and sets, inclusive), then the .csv file created by this
function has the suffix ``\_set\#stats.csv'', where ``\#'' is the set
number. If the ``set'' argument is ``allstats'', then the .csv file has
the suffix \_allstats.csv''. The output file contains the tables for all
cohorts with a non-zero population.

\hypertarget{summary.stats-summary.stats}{%
\section{\texorpdfstring{\texttt{summary.stats}:
summary.stats}{summary.stats: summary.stats}}\label{summary.stats-summary.stats}}

\hypertarget{description-53}{%
\subsection{Description}\label{description-53}}

Summary.stats constructs the table of exposure and dose statistics for
cohort entered into
\texttt{{[}summarize.chemical{]}(summarize.chemical)}.

\hypertarget{usage-53}{%
\subsection{Usage}\label{usage-53}}

\begin{Shaded}
\begin{Highlighting}[]
\CommentTok{\# S3 method for stats}
\FunctionTok{summary}\NormalTok{(x)}
\end{Highlighting}
\end{Shaded}

\hypertarget{arguments-38}{%
\subsection{Arguments}\label{arguments-38}}

\begin{itemize}
\tightlist
\item
  \texttt{x.}: Data set passed from summarize.chemical for a cohort
\end{itemize}

\hypertarget{details-30}{%
\subsection{Details}\label{details-30}}

Summary.stats is called by
\texttt{{[}summarize.chemical{]}(summarize.chemical)}. The input data
set ``y'' is one population cohort from the exposure data set passed
into \texttt{{[}summarize.chemical{]}(summarize.chemical)}.

\hypertarget{author-33}{%
\subsection{Author}\label{author-33}}

Kristin Isaacs, Graham Glen

\hypertarget{seealso-25}{%
\subsection{Seealso}\label{seealso-25}}

\texttt{{[}summarize.chemical{]}(summarize.chemical)}

\hypertarget{value-36}{%
\subsection{Value}\label{value-36}}

y A data frame object with 23 rows and 11 columns, with each column
being an exposure or dose variable, and each row containing a statistic
for that variable. For each expsoure varianle the total expsoure,
quantiles, mean, and SD.

\hypertarget{trimzero-trimzero}{%
\section{\texorpdfstring{\texttt{trimzero}:
Trimzero}{trimzero: Trimzero}}\label{trimzero-trimzero}}

\hypertarget{description-54}{%
\subsection{Description}\label{description-54}}

This function removes initial zeroes from CAS numbers.

\hypertarget{usage-54}{%
\subsection{Usage}\label{usage-54}}

\begin{Shaded}
\begin{Highlighting}[]
\FunctionTok{trimzero}\NormalTok{(x, y)}
\end{Highlighting}
\end{Shaded}

\hypertarget{arguments-39}{%
\subsection{Arguments}\label{arguments-39}}

\begin{itemize}
\tightlist
\item
  \texttt{x}: aCAS number. Defult is none
\item
  \texttt{y}: A dummy argument This helps with the removal of the zeros
  inteh CAS number
\end{itemize}

\hypertarget{details-31}{%
\subsection{Details}\label{details-31}}

Each CAS number has three parts, separated by underscores. The first
part is up to seven digits, but optionally, leading zeroes are omitted.
For example, formaldehyde may be either ``0000050\_00\_0'' or
``50\_00\_0''. SHEDS needs to match CAS numbers across input files, and
trimzero is used to ensure matching even when the input files follow
different conventions.

\hypertarget{author-34}{%
\subsection{Author}\label{author-34}}

Kristin Isaacs, Graham Glen

\hypertarget{value-37}{%
\subsection{Value}\label{value-37}}

y a shorter CAS number. If the initial part is all zero (as in
``0000000\_12\_3''), one zero is left in the first part (that is,
``0\_12\_3'' for this example).

\hypertarget{unpack-unpack}{%
\section{\texorpdfstring{\texttt{unpack}:
unpack}{unpack: unpack}}\label{unpack-unpack}}

\hypertarget{description-55}{%
\subsection{Description}\label{description-55}}

unpack

\hypertarget{usage-55}{%
\subsection{Usage}\label{usage-55}}

\begin{Shaded}
\begin{Highlighting}[]
\FunctionTok{unpack}\NormalTok{(}\AttributeTok{filelist =} \StringTok{""}\NormalTok{)}
\end{Highlighting}
\end{Shaded}

\hypertarget{details-32}{%
\subsection{Details}\label{details-32}}

This function is used when ShedsHT is run as an R package. Each time a
new working directory is chosen, use unpack() to convert the R data
objects into CSV files. Note that both /inputs and /output folders are
needed under the chosen directory.Unpack() may be used with a list of
object names, in which case just those objects are converted to CSV.
This is useful when re-loading one or more defaults into a folder where
some of the CSV files have changes, and should not be overwritten. A
blank argument or empty list means that all the csv and TXT files in the
R package are converted.

\hypertarget{author-35}{%
\subsection{Author}\label{author-35}}

Kristin Isaacs, Graham Glen

\hypertarget{update.specs-update.specs}{%
\section{\texorpdfstring{\texttt{update.specs}:
update.specs}{update.specs: update.specs}}\label{update.specs-update.specs}}

\hypertarget{description-56}{%
\subsection{Description}\label{description-56}}

``Specs'' is the list of run settings read from the run.file. ``Dt'' is
a data table of source-chemical combinations for which distributions
have been specified. ``Specs'' contains a list of chemicals to be
processed in the SHEDS run, and if any of these chemicals are missing
from the ``dt'' table, then update.specs removes them from the list.
Otherwise, ``specs'' is not altered.

\hypertarget{usage-56}{%
\subsection{Usage}\label{usage-56}}

\begin{Shaded}
\begin{Highlighting}[]
\CommentTok{\# S3 method for specs}
\FunctionTok{update}\NormalTok{(specs, dt)}
\end{Highlighting}
\end{Shaded}

\hypertarget{vpos-vpos}{%
\section{\texorpdfstring{\texttt{vpos}:
vpos}{vpos: vpos}}\label{vpos-vpos}}

\hypertarget{description-57}{%
\subsection{Description}\label{description-57}}

This function locates an item in a list.

\hypertarget{usage-57}{%
\subsection{Usage}\label{usage-57}}

\begin{Shaded}
\begin{Highlighting}[]
\FunctionTok{vpos}\NormalTok{(v, list)}
\end{Highlighting}
\end{Shaded}

\hypertarget{details-33}{%
\subsection{Details}\label{details-33}}

this one was written because it does not require additional R packages
and its behavior can be easily examined.

\hypertarget{author-36}{%
\subsection{Author}\label{author-36}}

Kristin Isaacs, Graham Glen

\hypertarget{write.persons-write.persons}{%
\section{\texorpdfstring{\texttt{write.persons}:
write.persons}{write.persons: write.persons}}\label{write.persons-write.persons}}

\hypertarget{description-58}{%
\subsection{Description}\label{description-58}}

This function writes demographic, exposure, and dose variables to a
separate output file for each chemical.

\hypertarget{usage-58}{%
\subsection{Usage}\label{usage-58}}

\begin{Shaded}
\begin{Highlighting}[]
\FunctionTok{write.persons}\NormalTok{(x, chem, set, specs)}
\end{Highlighting}
\end{Shaded}

\hypertarget{details-34}{%
\subsection{Details}\label{details-34}}

THis function rounds the variables to a reasonable precision so that the
files are more readable, as unrounded output is cluttered with entries
with (say) 14 digits, most of which are not significant.

\hypertarget{author-37}{%
\subsection{Author}\label{author-37}}

Kristin Isaacs, Graham Glen

\end{document}
